\documentclass[11pt]{article}
\usepackage[margin=0.9in]{geometry}
\usepackage{amsmath,amssymb,amsthm,mathtools}
\usepackage{booktabs,longtable,array}
\usepackage{enumitem}
\usepackage{microtype}
\usepackage[hidelinks]{hyperref}
\usepackage{xcolor}
\usepackage{enumitem}

\newtheorem{theorem}{Theorem}[section]
\newtheorem{proposition}[theorem]{Proposition}
\newtheorem{lemma}[theorem]{Lemma}
\newtheorem{corollary}[theorem]{Corollary}
\newtheorem{hypothesis}[theorem]{Hypothesis}
\newtheorem{claim}[theorem]{Claim}
\theoremstyle{definition}
\newtheorem{definition}[theorem]{Definition}
\newtheorem{status}[theorem]{Status}
\theoremstyle{remark}
\newtheorem{remark}[theorem]{Remark}

\newcommand{\GapCE}{\mathrm{Gap2CE}}
\newcommand{\SophieTI}{\mathrm{SOPHIE\mbox{-}FM\mbox{-}TYPEII}_{1/6}}
\newcommand{\FullFM}{\mathrm{FullFMTypeII}_{1/6}}
\newcommand{\FMRough}{\mathrm{FMRoughN2}}
\newcommand{\e}{\mathrm{e}}
\newcommand{\1}{\mathbf 1}
\newcommand{\statusbox}[1]{\begin{center}\fbox{\parbox{0.94\textwidth}{#1}}\end{center}}

\title{\textbf{Erd\H{o}s Problem \#287: A Standalone Verification Dossier}\\
\large Draft V12 -- explicit Ford certificate, smooth-parity counterguard, and corrected fixed-leakage frontier}
\author{Research draft -- not a solution claim}
\date{27 August 2026}

\begin{document}
\maketitle

\statusbox{\textbf{Purpose and status.}
This document is deliberately written so that the mathematical chain can be audited from this document alone.
It does not use ``see the Lean bank'', ``see the Gate-1A paper'', or ``see the Gate-0/2 archive'' as substitutes
for mathematical statements.  The finite Erd\H{o}s-287 reduction is proved explicitly below.  The Ford--Maynard
sequence conditions are written out explicitly.  Gate 1A and Gate 1B are reproduced as mathematical theorem
chains, with every load-bearing hypothesis stated.  Gate 2 is written as an explicit conditional endgame.
Where a statement remains unproved, the statement itself is displayed rather than hidden behind a status label.
No claim of Full Type II, twin-prime infinitude, or a solution of Problem \#287 is made.}

\begin{abstract}
Let $A$ be a finite set of distinct integers $1<n_1<\cdots<n_k$ with
$\sum_{a\in A}1/a=1$.  Erd\H{o}s Problem \#287 asks whether some consecutive gap must be at least three.
We first give a self-contained finite reduction.  The reciprocal identity implies a top-layer $p$-adic
congruence.  Together with the elementary half-range bound $N\le\lfloor M/2\rfloor$, this shows that a prime
$q>3$ with $3q>M$ and $q^2>M$ forces both $q$ and $2q$ to be absent, while any prime $p\in(M/2,M]$ is absent.
Hence an interior prime pair $p=2q\pm1$ produces two adjacent holes and contradicts the gap-$\le2$ hypothesis.
A fixed band $q\in(3X/4,4X/5)$ with $X=\lfloor M/2\rfloor$ suffices for every $M\ge9$.

The analytic task is therefore to force such a Sophie-type witness.  We define the comparison density
\[
 B(n)=\mathfrak S_2\prod_{\ell\mid n,\ell>2}\frac{\ell-1}{\ell-2},
\]
and prove $Y^{-1}\sum_{m\le Y}B(dm)\to d/\varphi(2d)$, which is the exact local mean required by the two
affine forms $2n\pm1$.  We then write out the Ford--Maynard Type-I, Type-II, growth, comparison and bounded-class hypotheses,
the $1/6$ parameter geometry, the Gate-0 specialization, the complete current Gate-1A conditional mechanism,
the complete current Gate-1B conductor reduction and its late pair-modulus reorganisation, the source-exhaustive
Full-Type-II socket, and the Gate-2 rough endgame.  Draft V12 preserves the corrected Ford--Maynard $N_2$ replay
from Draft V8 and incorporates the 26/27 August Gate-1B and Ford-side updates: H7 short-short now has a conditional joint-prime compiler
with relative capacity $Y^{-1/2}=X^{-1/18}$; the H8 analogue has the same research-level fixed-power capacity but
still lacks its formal source/robustness pin; the QK56 full-covariance parent is promoted internally to an analytic pass modulo literal source pins,
while the principal Gate-1B analytic candidate moves to the balanced shifted-quotient child after a mandatory
literal $TT^*$ source recovery.  The canonical Gate-1B additive zero mode remains an exact algebraic pass.
Downstream, Ford--Maynard five-collar geometry now gives a source-specific organisation of the rough $N_2$
sector, the grouped equation (7.23) remains the proof-specific Type-II landing point, and the first unresolved
Gate-to-Ford interface is the literal weighted packet dictionary labelled SOURCE-ADAPTER45.
A subsequent hostile census corrects the first interpretation of this fixed-certificate pivot.  Ford--Maynard
supply an explicit certificate $g_0$ at their central point; after a small fixed shrink one obtains a literal
certificate $g_*$ and arithmetic weight $H_*$.  The fixed-certificate transference identity is valid, but the
resulting leakage family does \emph{not} collapse to only H8/H9 and QK56 packets.  Already the $k=0$, $J=\varnothing$
cell contains an exact smooth-parity packet with the truncated M\"obius divisor detector
$\sum_{d\mid n,\,d\le n^{1/2-\varepsilon_*}}\mu(d)$.  Balanced squarefree cells show nonzero contributions at
orders $7,8,9,10,11,12,\ldots$, so an ``H8/H9-only'' census is false.  The generic seven-prime quadratic-Kummer
estimate remains a useful conditional child provider with relative capacity $X^{-1/18+o(1)}$, but it is no longer
the controlling first wall.  The first literal analytic open in the fixed-certificate route is the source-exact
smooth-parity packet \texttt{287-FIXED-CERTIFICATE-SMOOTH-PARITY45}; the full fixed-certificate leakage estimate
and Problem~\#287 remain open.
\end{abstract}

\tableofcontents
\clearpage

\part{Erd\H{o}s \#287: the finite reduction written from first principles}

\section{Problem, notation, and the hypothetical counterexample}

Let
\[
 A=\{n_1<\cdots<n_k\}\subset\mathbb N,\qquad n_1>1,
 \qquad \sum_{i=1}^k\frac1{n_i}=1.
\]
Write
\[
 N:=n_1,\qquad M:=n_k.
\]
Assume for contradiction that
\[
 n_{i+1}-n_i\le2\qquad(1\le i<k).
\tag{G2}
\]
Call an integer $h\in[N,M]$ a \emph{hole} if $h\notin A$.

\begin{lemma}[Two adjacent holes are impossible]\label{lem:isolated}
Under \emph{(G2)}, there do not exist $x,x+1\in[N,M]$ which are both holes.
\end{lemma}
\begin{proof}
Since $N,M\in A$, neither endpoint can be a hole.  Thus if $x,x+1$ are holes there is a largest denominator
$a\in A$ with $a<x$ and a smallest denominator $b\in A$ with $b>x+1$.  Then $b-a\ge3$, contradicting (G2).
\end{proof}

\section{The half-range placement bound}

\begin{lemma}[Half-range placement]\label{lem:half}
Every hypothetical counterexample satisfies
\[
 N\le \left\lfloor\frac M2\right\rfloor.
\]
\end{lemma}
\begin{proof}
Suppose instead $N\ge\lfloor M/2\rfloor+1$.  Then $2N\ge M+1$, hence
\[
 k\le M-N+1\le N.
\]
Because the denominators are distinct and at least $N$,
\[
 \sum_{i=1}^k\frac1{n_i}
 \le \frac1N+\frac{k-1}{N+1}
 < \frac{k}{N}\le1,
\]
contradicting the reciprocal identity.  (If $k=1$, the sole denominator would be $1$, excluded by hypothesis.)
\end{proof}

\section{The top-layer congruence derived directly from the reciprocal identity}

\begin{theorem}[Top-layer congruence]\label{thm:TL}
Fix a prime $p$.  Let
\[
 E:=\max_{a\in A}v_p(a),
 \qquad I_p:=\{a\in A:v_p(a)=E\}.
\]
For each $a\in I_p$ write $a=p^E m_a$ with $p\nmid m_a$.  Then
\[
 \boxed{\sum_{a\in I_p}m_a^{-1}\equiv0\pmod p.}
\tag{TL}
\]
Here $m_a^{-1}$ is the inverse of $m_a$ modulo $p$.
\end{theorem}
\begin{proof}
Let $L=\operatorname{lcm}(A)$ and write $L=p^E L_0$ with $p\nmid L_0$.  Multiplying
$\sum_{a\in A}1/a=1$ by $L$ gives
\[
 L=\sum_{a\in A}\frac La.
\]
Reduce modulo $p$.  If $v_p(a)<E$, then $L/a$ is divisible by $p$, so these terms vanish.  If
$a=p^E m_a\in I_p$, then
\[
 \frac La=\frac{L_0}{m_a}\equiv L_0m_a^{-1}\pmod p.
\]
Since $L\equiv0\pmod p$ and $L_0$ is a unit modulo $p$,
\[
 0\equiv L_0\sum_{a\in I_p}m_a^{-1}\pmod p,
\]
which proves (TL).
\end{proof}

\section{Immediate prime exclusions}

\begin{lemma}[A two-multiple prime fibre is empty]\label{lem:q2}
Let $q>3$ be prime and suppose
\[
 3q>M,\qquad q^2>M.
\]
Then neither $q$ nor $2q$ belongs to $A$.
\end{lemma}
\begin{proof}
The only multiples of $q$ in $[1,M]$ are $q$ and $2q$, and $q^2>M$ ensures both have $q$-adic valuation one.
If at least one belonged to $A$, the maximal $q$-adic fibre would be one of
\[
 \{q\},\qquad\{2q\},\qquad\{q,2q\}.
\]
By Theorem~\ref{thm:TL}, these would force respectively
\[
 1\equiv0,\qquad 2^{-1}\equiv0,\qquad1+2^{-1}=\frac32\equiv0\pmod q,
\]
all impossible for $q>3$.
\end{proof}

\begin{lemma}[Every top-half prime is a hole]\label{lem:tophalf}
If $p$ is prime and
\[
 \frac M2<p\le M,
\]
then $p\notin A$.
\end{lemma}
\begin{proof}
The only multiple of $p$ in $[1,M]$ is $p$.  If $p\in A$, the top $p$-adic fibre is the singleton $\{p\}$,
and (TL) gives $1\equiv0\pmod p$.
\end{proof}

\section{The direct $\pm$ Sophie blockers}

\begin{proposition}[Plus-Sophie blocker]\label{prop:plus}
Let $q>3$ and $p=2q+1$ be primes.  If
\[
 p\le M<3q,
 \qquad q^2>M,
\]
then the assumed counterexample cannot exist.
\end{proposition}
\begin{proof}
Lemma~\ref{lem:q2} makes $2q=p-1$ a hole.  Since
\[
 M<3q<4q+2=2p,
\]
we have $p>M/2$, so Lemma~\ref{lem:tophalf} makes $p$ a hole.  Lemma~\ref{lem:half} gives
$N\le\lfloor M/2\rfloor<p-1$, and $p\le M$.  Hence $p-1,p$ are adjacent holes in $[N,M]$, contradicting
Lemma~\ref{lem:isolated}.
\end{proof}

\begin{proposition}[Minus-Sophie blocker]\label{prop:minus}
Let $q>3$ and $p=2q-1$ be primes.  If
\[
 p+1\le M<3q,
 \qquad q^2>M,
\]
then the assumed counterexample cannot exist.
\end{proposition}
\begin{proof}
Lemma~\ref{lem:q2} makes $2q=p+1$ a hole.  Since $q>3$,
\[
 M<3q<4q-2=2p,
\]
so $p>M/2$ and Lemma~\ref{lem:tophalf} makes $p$ a hole.  Lemma~\ref{lem:half} gives
$N\le\lfloor M/2\rfloor<p$, while $p+1\le M$.  Hence $p,p+1$ are adjacent holes, contradiction.
\end{proof}

\section{A fixed analytic band}

Put
\[
 X:=\left\lfloor\frac M2\right\rfloor
\]
and fix a non-negative $W\in C_c^\infty((3/4,4/5))$.

\begin{lemma}[Band arithmetic]\label{lem:band}
For every integer $M\ge9$, if $W(q/X)>0$, then
\[
 \boxed{q>3,\qquad M<3q,\qquad 2q+1\le M,\qquad q^2>M.}
\]
\end{lemma}
\begin{proof}
We have $3X/4<q<4X/5$ and $X=\lfloor M/2\rfloor$.  Since $M\ge9$, $X\ge4$, hence
$q>3X/4\ge3$, and the strict inequality gives $q>3$.

For the lower Sophie-band inequality, $X\ge(M-1)/2$, so
\[
 3q>\frac{9X}{4}\ge\frac{9(M-1)}8\ge M,
\]
where the last inequality is equivalent to $M\ge9$.
Also $q<4X/5\le2M/5$, hence
\[
 2q+1<\frac{4M}{5}+1\le M\qquad(M\ge5).
\]
Finally, strictness is important at the endpoint.  From
\[
 q>\frac{3X}{4}\ge\frac{3(M-1)}8
\]
and
\[
 \left(\frac{3(M-1)}8\right)^2-M
 =\frac{9M^2-82M+9}{64}
 =\frac{9(M-9)(M-1/9)}{64}\ge0
\]
for $M\ge9$, we obtain
\[
 q^2>\left(\frac{3(M-1)}8\right)^2\ge M.
\]
Thus the threshold $M\ge9$ is valid; no replacement by $M\ge10$ is required.
\end{proof}

\begin{corollary}[Finite landing theorem]\label{cor:landing}
For $M\ge9$, if there exists a prime $q$ with $W(q/X)>0$ such that at least one of $2q-1,2q+1$ is prime,
then no counterexample with maximum denominator $M$ exists.
\end{corollary}

\section{The exact comparison density for the two affine prime forms}

Define
\[
 \mathfrak S_2:=\prod_{\ell>2}\frac{\ell(\ell-2)}{(\ell-1)^2},
 \qquad
 B(n):=\mathfrak S_2\prod_{\substack{\ell\mid n\\\ell>2}}\frac{\ell-1}{\ell-2}.
\]
Let
\[
 a_X(n):=W(n/X)\bigl(\Lambda(2n-1)+\Lambda(2n+1)\bigr),
\]
\[
 b_X(n):=4W(n/X)B(n),\qquad w_X(n):=a_X(n)-b_X(n).
\]
For odd squarefree $e$ set
\[
 g(e):=\prod_{\ell\mid e}\frac1{\ell-2},
\]
and put $g(e)=0$ otherwise.  Then
\[
 B(n)=\mathfrak S_2\sum_{e\mid n}g(e).
\tag{7.1}
\]

\begin{proposition}[Local mean]\label{prop:Bmean}
For each fixed positive integer $d$,
\[
 \lim_{Y\to\infty}\frac1Y\sum_{m\le Y}B(dm)=\frac d{\varphi(2d)}.
\]
\end{proposition}
\begin{proof}
The series
\[
 \sum_e\frac{|g(e)|}{e}=\prod_{\ell>2}\left(1+\frac1{\ell(\ell-2)}\right)
\]
converges absolutely.  Truncate at $e\le Z$, average, and then let $Y\to\infty$ followed by $Z\to\infty$.
Since the density of $m$ with $e\mid dm$ is $(e,d)/e$,
\[
 \lim_{Y\to\infty}\frac1Y\sum_{m\le Y}B(dm)
 =\mathfrak S_2\sum_e g(e)\frac{(e,d)}e.
\]
For an odd prime $\ell\nmid d$, the Euler factor is
\[
 \frac{\ell(\ell-2)}{(\ell-1)^2}\left(1+\frac1{\ell(\ell-2)}\right)=1.
\]
For $\ell\mid d$, it is
\[
 \frac{\ell(\ell-2)}{(\ell-1)^2}\left(1+\frac1{\ell-2}\right)=\frac\ell{\ell-1}.
\]
The prime $2$ contributes no factor.  Hence the product is
$\prod_{\ell\mid d,\ell>2}\ell/(\ell-1)=d/\varphi(2d)$.
\end{proof}

\begin{remark}[Why the factor $4$ is exact]
For each sign, the affine form $2dm\pm1$ has von Mangoldt mean $2d/\varphi(2d)$, so the two signs together
have mean $4d/\varphi(2d)$, exactly matching $4B(dm)$.
\end{remark}

\part{The Ford--Maynard framework and Gate 0, written explicitly}

\section{The exact Ford--Maynard sequence class}

This section writes out the external sieve framework used downstream.  The notation follows Ford--Maynard
\cite{FM24}, but the formulas needed here are reproduced so that the application can be audited locally.
Fix parameters
\[
 0<\gamma<1,\qquad0\le\theta<\frac12,\qquad0<\nu\le1-\theta,
\]
and non-negative sequences $(a_n),(b_n)$ supported on $x/2<n\le x$.  Put $w_n=a_n-b_n$.

\subsection{Type I, Type II, and the growth condition}

The Type-I hypothesis is
\[
 \boxed{
 \sum_{m\le x^\gamma}\tau^B(m)
 \max_{I\ \mathrm{interval}}
 \left|\sum_{\substack{x/2<mn\le x\\ n\in I}}w_{mn}\right|
 \le \frac{x}{(\log x)^B}.}
\tag{I}
\]
The Type-II hypothesis is universal in the coefficients:
\[
 \boxed{
 \left|\sum_{\substack{(x/2)^\theta<m\le x^{\theta+\nu}\\x/2<mn\le x}}
 \xi_m\kappa_n w_{mn}\right|
 \le\frac{x}{(\log x)^B}}
\tag{II}
\]
for every pair of complex sequences satisfying
$|\xi_m|\le\tau^B(m)$ and $|\kappa_n|\le\tau^B(n)$.
The lower endpoint is $(x/2)^\theta$, not $x^\theta$.

The ordinary growth condition is
\[
 \boxed{\sum_n|w_n|\tau(n)\le x(\log x)^\varpi,
 \qquad w_n\ge -x^{\nu/10}\quad(x/2<n\le x).}
\tag{w}
\]

\subsection{The comparison-sequence conditions}

Let $P=(\gamma,\theta,\nu)$.  Define $R(P)$ to be the set of finite vectors
$u=(u_1,\ldots,u_k)$ whose components lie in $(0,1-\gamma)$, sum to $1$, and have no proper subsum in
$[\theta,\theta+\nu]$.  Let $C(R(P))$ denote all coagulations of vectors in $R(P)$.

The first comparison condition is the prime-mass lower bound
\[
 \boxed{\sum_p b_p\ge \frac{x}{(\log x)^\varpi}.}
\tag{b.1}
\]
The second is a generalized prime-number-theorem condition.  For every
$2\le k\le1/\nu$, every convex
\[
 T\subset\{u\in\mathbb R^k\cap C(R(P)):\nu\le u_1\le\cdots\le u_k\},
\]
and every $1$-bounded Lipschitz function $f$ of Lipschitz constant at most $1$,
\[
 \boxed{
 \sum_{\substack{n=p_1\cdots p_k\\p_1\le\cdots\le p_k}}
 b_n f\!\left(\frac{\log p_1}{\log n},\ldots,\frac{\log p_k}{\log n}\right)
 =\left(\sum_p b_p\right)
 \left(\int_T\frac{f(u)\,du}{u_1\cdots u_k}+E\right),
 \quad |E|\le\frac1B.}
\tag{b.2}
\]

Ford--Maynard's ordinary sequence class consists of non-negative $(a_n),(b_n)$ satisfying
(b.1), (b.2), with $w=a-b$ satisfying (I), (II), and (w).

\subsection{The bounded class}

For the bounded constants $C^\pm_{\mathrm{bd}}$, one additionally imposes, for some fixed $\varrho\ge1$,
\[
 \boxed{|w_n|\le\tau(n)^\varrho\quad(x/2<n\le x),
 \qquad\sum_p b_p\ge\frac{x}{\varrho\log x}.}
\tag{4.1}
\]
This extra pointwise condition is the source of the $N_2$ mismatch discussed later.

\section{The published Ford--Maynard $1/6$ region and the class/shrinkage firewall}

\begin{theorem}[Ford--Maynard, special family]\label{thm:FM27}
Fix $\varrho\ge1$.  For
\[
 P=(1/2,0,\nu),\qquad P_\varepsilon=(1/2-\varepsilon,\varepsilon,\nu-2\varepsilon),
\]
Ford--Maynard prove that if $\nu\ge0.1663$, then
\[
 C^-(1/2,0,\nu)>0,
 \qquad
 \lim_{\varepsilon\to0^+}C^-_{\rm bd}(P_\varepsilon;\varrho)>0.
\]
\end{theorem}

For the operating value $\nu_0=1/6$,
\[
 \frac16-0.1663=\frac{11}{30000}>0.
\]
The theorem therefore supplies two \emph{different} published facts:
\begin{enumerate}[label=(\roman*)]
\item positivity for the \emph{ordinary} class at the unshrunk point $P=(1/2,0,1/6)$;
\item epsilon-stable positivity for the \emph{bounded} class at the shrunk points
\[
 P_\varepsilon=\left(\frac12-\varepsilon,\ \varepsilon,\ \frac16-2\varepsilon\right).
\]
\end{enumerate}
It does \textbf{not} by itself assert ordinary-class positivity at $P_\varepsilon$.  Since the 287 prime source
contains von Mangoldt weights, the pointwise bounded-class condition $|w_n|\le\tau(n)^\varrho$ is unavailable.
Draft V8 therefore never applies $C^-_{\rm bd}(P_\varepsilon;\varrho)$ directly to the 287 source.

The exact Type-II interval at the shrunk point is
\[
 \boxed{(x/2)^\varepsilon<m\le x^{1/6-\varepsilon}.}
\tag{FM-window}
\]
The role of the published bounded theorem is instead to identify the central Ford--Maynard geometry and a
positive limiting margin.  The source-specific application below replays the proof of their Theorem~8.2 and
replaces the only pointwise-bounded exceptional step by an explicit $N_2$ estimate.

\begin{status}[Class/shrinkage firewall]
There is no theorem in this dossier asserting
\[
 C^-(P_\varepsilon)>0
\]
for the ordinary class.  Such a theorem is unnecessary if the source-specific replay in
Theorem~\ref{thm:Gate2} is justified.  Conversely, if the replay/localization failed, ordinary-class positivity
at a shrunk point would become a genuinely new obligation.
\end{status}

\section{Gate 0 before specialization: the actual shifted-prime source}
\label{sec:g0-actual}

The original programme does not begin with the 287 source $2n\pm1$.  Its Gate-0 physical sequence is a
single fixed shift.  We record it because a reader must be able to distinguish the programme theorem from the
later 287 specialization.

For $z\ge2$ define the odd primorial and its finite sieve density
\[
 P_0(z):=\prod_{\substack{p\le z\\p\ {
m prime},\ p>2}}p,
 \qquad
 V_0(z):=\frac12\prod_{\substack{p\le z\\p\ {
m prime},\ p>2}}\left(1-\frac1p\right).
\tag{G0-1}
\]
Let $I_x(n)=\1_{\{x/2<n\le x\}}$.  The shifted-prime detector, residue-aware comparison sequence and centered
difference are
\[
 a_x(n):=I_x(n)\frac{\Lambda(n+2)}{\log x},
\tag{G0-2}
\]
\[
 b_{x,z}(n):=
 \frac{I_x(n)\,\1_{\{n\ {
m odd},\ (n+2,P_0(z))=1\}}}{\log x\,V_0(z)},
 \qquad
 w_{x,z}(n):=a_x(n)-b_{x,z}(n).
\tag{G0-3}
\]
Both $a_x$ and $b_{x,z}$ are non-negative and supported on $(x/2,x]$; moreover $b_{x,z}(n)=0$ for even $n$.
These are direct consequences of the definitions.

\subsection{Finite local densities and the twin Euler product}

Define
\[
 D_z(m):=\prod_{\substack{p\le z,\ p>2\\p\mid m}}
 \left(1-\frac1p\right)^{-1}.
\tag{G0-4}
\]
Then the exact finite local-density identity is
\[
 \boxed{
 \frac12\prod_{\substack{p\le z,\ p>2\\p\nmid m}}
 \left(1-\frac1p\right)=V_0(z)D_z(m).}
\tag{G0-5}
\]
Indeed split the full product defining $V_0(z)$ into primes dividing $m$ and primes not dividing $m$, and cancel
the first group against the inverse factors in $D_z(m)$.

Consequently
\[
 D_z(1)=1,
\tag{G0-6}
\]
and for an odd prime $p$,
\[
 D_z(p)=
 \begin{cases}
 (1-1/p)^{-1},&p\le z,\\
 1,&p>z.
 \end{cases}
\tag{G0-7}
\]
For every $a\ge1$, $D_z(p^a)=D_z(p)$; more generally $D_z(m)$ depends only on the radical of odd $m$.

A second finite product is
\[
 W_0(z):=\prod_{\substack{p\le z\\p>2}}
 \left(1-\frac1{p-1}\right).
\tag{G0-8}
\]
The elementary local identity
\[
 \frac{p(p-2)}{(p-1)^2}=1-\frac1{(p-1)^2}
\tag{G0-9}
\]
gives
\[
 \boxed{
 \frac{W_0(z)}{V_0(z)}
 =2\prod_{\substack{p\le z\\p>2}}
 \frac{p(p-2)}{(p-1)^2}
 =2\prod_{\substack{p\le z\\p>2}}
 \left(1-\frac1{(p-1)^2}\right).}
\tag{G0-10}
\]
This is an exact \emph{finite} singular-series identity; no convergence assertion is needed for it.

\subsection{The elementary pointwise comparison bound}

If a Mertens lower bound
\[
 V_0(z)\ge\frac{c}{\log z},\qquad c>0,
\tag{G0-11}
\]
is available, then directly from (G0-3)
\[
 0\le b_{x,z}(n)\le\frac{\log z}{c\log x}.
\tag{G0-12}
\]
Thus pointwise boundedness of the comparison sequence is an ordinary finite-sieve consequence once (G0-11) is
pinned; it is not a Type-II input.

\section{Gate 0 / Gate 1 front algebra: centering, Fourier separation, CRT and the fixed determinant}
\label{sec:g01-front}

This section writes out the exact finite identities which precede the later analytic estimates.  They are useful
because they show precisely where the expected main term disappears and where the fixed shift $2$ enters.

\subsection{Nonzero additive orthogonality and E-separation}

Write $e_q(t)=e^{2\pi i t/q}$.  For $q>0$,
\[
 \sum_{a\bmod q}e_q(at)=q\,\1_{q\mid t},
\tag{G01-1}
\]
so after deleting the zero frequency,
\[
 \boxed{
 \sum_{\substack{a\bmod q\\a\ne0}}e_q(at)
 =q\,\1_{q\mid t}-1,}
\tag{G01-2}
\]
and hence
\[
 \frac1q\sum_{a\ne0}e_q(at)=\1_{q\mid t}-\frac1q.
\tag{G01-3}
\]
For a finite set $S$ and coefficients $c_x$, put
\[
 C_q:=\sum_{\substack{x\in S\\q\mid x+2}}c_x,
 \qquad C:=\sum_{x\in S}c_x.
\]
Using (G01-3),
\[
 \boxed{
 C_q=\frac Cq+
 \frac1q\sum_{a\ne0}e_q(2a)\sum_{x\in S}c_xe_q(ax).}
\tag{G01-4}
\]
Therefore for \emph{any} proposed expected value $E$,
\[
 C_q-E=\left(\frac Cq-E\right)
 +\frac1q\sum_{a\ne0}e_q(2a)\widehat c_q(a).
\tag{G01-5}
\]
The nonzero-frequency term is independent of the choice of $E$.  Thus an incorrect source normalization cannot
be repaired by changing the nonzero Fourier analysis; the zero mode and the nonzero modes are logically separate.

If $n\bar n\equiv1\pmod q$, then
\[
 \boxed{q\mid m+2\bar n\quad\Longleftrightarrow\quad q\mid mn+2.}
\tag{G01-6}
\]
This follows by multiplying the first congruence by $n$ and using $n\bar n\equiv1$.

\subsection{CRT centering}

For $r\ge1$ define
\[
 \rho_r(y):=\1_{r\mid y}-\frac1r.
\]
If $(d,p)=1$, then $\1_{dp\mid y}=\1_{d\mid y}\1_{p\mid y}$, whence
\[
 \boxed{
 \rho_{dp}(y)=\rho_d(y)\rho_p(y)+\frac{\rho_d(y)}p+\frac{\rho_p(y)}d.}
\tag{G01-7}
\]
Thus a composite centered defect has a double-centered piece and two one-coordinate pieces; none may be silently
removed.

At the frequency level,
\[
 \frac{a_d}{d}+\frac{a_p}{p}=\frac{a_dp+a_pd}{dp},
\tag{G01-8}
\]
and therefore
\[
 \boxed{e_{dp}((a_dp+a_pd)y)=e_d(a_dy)e_p(a_py).}
\tag{G01-9}
\]
This is the exact CRT frequency split underlying the four mode classes: double-centered, $d$-only, $p$-only,
and zero mode.

\subsection{Shift inverses and the fixed-shift determinant}

Suppose
\[
 v=n+hq,
 \qquad q\bar q\equiv1\pmod n,
 \qquad v\bar v\equiv1\pmod n.
\]
Then
\[
 \boxed{\bar q\equiv h\bar v\pmod n,}
\tag{G01-10}
\]
so for every integer $a$,
\[
 e_n(-2a\bar q)=e_n(-2ah\bar v).
\tag{G01-11}
\]
The proof of (G01-10) is direct: substitute $v=n+hq$ and compare the two inverse identities modulo $n$.

If in addition
\[
 mn+2=qr_0,
 \qquad v=n+hq,
 \qquad s=r_0+hm,
\]
then
\[
 \boxed{r_0v-ns=2h.}
\tag{G01-12}
\]
Indeed
$r_0(n+hq)-n(r_0+hm)=h(qr_0-mn)=2h$.
The determinant shift is therefore literally the original fixed shift $2$, multiplied only by the translation
parameter $h$.

\subsection{Prime covariance and the exact second moment}

For a finite prime set $\mathcal P$, define
\[
 K_{\mathcal P}(u,v):=\sum_{p\in\mathcal P}\rho_p(u)\rho_p(v).
\tag{G01-13}
\]
Expanding gives
\[
 \boxed{
 K_{\mathcal P}(u,v)
 =\sum_{\substack{p\in\mathcal P\\p\mid u,\ p\mid v}}1
 -\sum_{\substack{p\in\mathcal P\\p\mid u}}\frac1p
 -\sum_{\substack{p\in\mathcal P\\p\mid v}}\frac1p
 +\sum_{p\in\mathcal P}\frac1{p^2}.}
\tag{G01-14}
\]
On the diagonal,
\[
 K_{\mathcal P}(Y,Y)=
 \sum_{\substack{p\in\mathcal P\\p\mid Y}}\left(1-\frac2p\right)
 +\sum_{p\in\mathcal P}\frac1{p^2}.
\tag{G01-15}
\]
For amplitudes $A_i$ and integers $Y_i$, the centered form
\[
 F_p:=\sum_iA_i\rho_p(Y_i)
\]
satisfies the exact Hermitian identity
\[
 \boxed{
 \sum_{p\in\mathcal P}|F_p|^2
 =\sum_{i,j}A_i\overline{A_j}\,K_{\mathcal P}(Y_i,Y_j).}
\tag{G01-16}
\]
No analytic saving is contained in (G01-16); it only exposes the diagonal and off-diagonal kernels exactly.

\subsection{Direct Gauss reassembly}

For a unit $A$ modulo $c$, the congruence $Ax\equiv B\pmod c$ has the unique solution $x\equiv A^{-1}B$.
Hence
\[
 \boxed{
 \sum_{\substack{x\bmod c\\Ax\equiv B\ (c)}}e_c(x)=e_c(A^{-1}B).}
\tag{G01-17}
\]
In the direct family this produces the physical phase
\[
 e_c\!\left(2kN\bigl((m+kr)p_1p_2\bigr)^{-1}\right)
\tag{G01-18}
\]
on the unit sector.  If $A$ is not a unit, the correct solvability criterion is
\[
 \boxed{\exists x:\ Ax\equiv B\pmod c\quad\Longleftrightarrow\quad (A,c)\mid B.}
\tag{G01-19}
\]
Thus non-unit strata have an explicit multiplicity/gcd structure and cannot be folded into the unique-solution
formula.

\subsection{Nine-block regrouping and the completion exponents}

For nine equal logarithmic blocks, each has mass $1/9$.  The balanced binary regrouping has exponents
\[
 U=X^{4/9+o(1)},\qquad V=X^{5/9+o(1)},\qquad UV=X^{1+o(1)}.
\tag{G01-20}
\]
At the hard modulus exponent
\[
 Q=X^{13/18+o(1)},\qquad \sqrt Q=X^{13/36+o(1)},
\tag{G01-21}
\]
one has
\[
 \frac49-\frac{13}{36}=\frac1{12},
 \qquad
 \frac59-\frac{13}{36}=\frac7{36},
\tag{G01-22}
\]
so both regrouped variables lie above the square-root modulus scale.  The natural progression-count exponent is
\[
 \frac49+\frac59-\frac{13}{18}=\frac5{18},
\tag{G01-23}
\]
and
\[
 \frac{13}{36}-\frac5{18}=\frac1{12}.
\tag{G01-24}
\]
These are exponent identities only; they do not themselves prove a square-root error term or an analytic
completion theorem.


\section{Gate 0: the 287 sequence dictionary}

We now use the 287 sequences
\[
 a_X(n)=W(n/X)(\Lambda(2n-1)+\Lambda(2n+1)),\quad
 b_X(n)=4W(n/X)B(n),\quad w_X=a_X-b_X.
\]
Because $\operatorname{supp}W\subset(3/4,4/5)$, these sequences are supported inside $(X/2,X]$.

\subsection{The comparison prime mass (b.1) and the bounded-class prime mass}

For prime $p$ in the support,
\[
 B(p)=\mathfrak S_2\frac{p-1}{p-2}=\mathfrak S_2(1+O(1/p)).
\]
The smooth prime number theorem gives
\[
 \sum_p b_X(p)
 =4\mathfrak S_2\sum_p W(p/X)+o(X/\log X)
 \sim 4\mathfrak S_2\left(\int_0^\infty W(t)\,dt\right)\frac X{\log X}.
\tag{G0-b1}
\]
Therefore (b.1) holds for every fixed $\varpi>1$ once $X$ is large, and the second inequality in (4.1) holds
for every fixed sufficiently large $\varrho$.

\subsection{The ordinary growth condition (w): an explicit logarithmic proof}

The comparison density has a fixed logarithmic bound, not merely an $X^{o(1)}$ bound.
For every odd prime $\ell$,
\[
 \frac{\ell-1}{\ell-2}
 =\frac{\ell}{\ell-1}
 \left(1+\frac{1}{\ell(\ell-2)}\right).
\]
Since
\[
 \prod_{\ell>2}\left(1+\frac{1}{\ell(\ell-2)}\right)<\infty,
\]
we obtain
\[
 B(n)\ll \prod_{\substack{\ell\mid n\\\ell>2}}\frac{\ell}{\ell-1}
 \ll \log\log(3n),
\tag{G0-Bpoint}
\]
where the last estimate is the standard maximal-order/Mertens bound for $n/\varphi(n)$.

On the support $n\asymp X$,
\[
 a_X(n)\ll_W\log X,\qquad b_X(n)\ll_W\log\log X,
\]
so
\[
 |w_X(n)|\ll_W\log X.
\]
Using $\sum_{n\asymp X}\tau(n)\ll X\log X$ gives the explicit fixed-power logarithmic estimate
\[
 \boxed{\sum_n|w_X(n)|\tau(n)\ll_W X(\log X)^2.}
\tag{G0-w1}
\]
Thus the first part of (w) holds, for example, with any fixed $\varpi>2$.
Moreover $a_X(n)\ge0$ and
\[
 w_X(n)\ge-b_X(n)\gg_W-\log\log X\ge-X^{\nu/10}
\]
for every fixed $\nu>0$ and sufficiently large $X$.  Hence the complete ordinary condition (w) is an
\textbf{elementary PASS} for the 287 sequence.

\subsection{Comparison condition (b.2)}

For $n=p_1\cdots p_k$ in the region occurring in (b.2), every $p_i\ge n^\nu$, hence $k\le1/\nu$ and
\[
 B(n)=\mathfrak S_2\prod_{i=1}^k\frac{p_i-1}{p_i-2}
 =\mathfrak S_2(1+O_\nu(X^{-\nu+o(1)})).
\]
Consequently (b.2) reduces to the classical fixed-$k$ multivariate prime number theorem with a smooth radial
weight.  Repeated prime-number-theorem partial summation gives, for each fixed admissible $T$ and Lipschitz $f$,
\[
 \sum_{\substack{n=p_1\cdots p_k\\p_1\le\cdots\le p_k}}
 b_X(n)f(v(n))
 =\left(\sum_p b_X(p)\right)
 \left(\int_T\frac{f(u)\,du}{u_1\cdots u_k}+o(1)\right).
\tag{G0-b2}
\]
For the Ford--Maynard definition one must make the $o(1)$ uniform over their normalized class of convex sets
and Lipschitz functions; this is a standard compactness/partial-summation task but should be written with an
explicit reference or proof before publication.

\subsection{Type I from a weighted maximal Bombieri--Vinogradov theorem}

The prime part of the Type-I sum is
\[
 \sum_{m\le X^{1/2-\varepsilon}}\tau^A(m)
 \max_I\left|
 \sum_{\substack{X/2<mn\le X\\n\in I}}
 W(mn/X)\Lambda(2mn\pm1)-\text{expected main term}\right|.
\]
For fixed $m$, the affine prime lies in a residue class $\pm1\pmod{2m}$.  Since
$2m\le2X^{1/2-\varepsilon}$, a weighted maximal Bombieri--Vinogradov theorem at level
$X^{1/2}(\log X)^{-C}$ has a fixed-power margin.  Partial summation inserts $W$, and divisor-weighted versions
absorb $\tau^A(m)$.  The comparison term $4B(mn)$ is handled by (7.1) and Proposition~\ref{prop:Bmean} in
progressions.  Thus the exact Gate-0 theorem needed is the following.

\begin{hypothesis}[287 Type-I source theorem]\label{hyp:G0I}
For every $A>0$ and fixed $\varepsilon>0$,
\[
 \boxed{
 \sum_{m\le X^{1/2-\varepsilon}}\tau^A(m)
 \max_I\left|\sum_{\substack{X/2<mn\le X\\n\in I}}w_X(mn)\right|
 \ll_{A,W,\varepsilon}X(\log X)^{-A}.}
\tag{287-I}
\]
\end{hypothesis}

The preceding calculation explains why this is a source-dictionary theorem at the classical $1/2$ level,
not the parity-breaking frontier.  The literal weighted/maximal citation and the comparison-progression error
remain to be pinned.

\part{Gate 1A: the complete current conditional mechanism}

\section{Gate-1A parameters and canonical physical source}

Let the Gate-1A ambient parameter be $X_1\to\infty$ and set
\[
 M=X_1^{1/3},\qquad R=X_1^a,\qquad L=X_1^b,
\]
\[
 H=X_1^{a+2b-2/3},\qquad K=M/R,\qquad U=L/H,\qquad D=L^2/H.
\]
Then
\[
 RK=M,\qquad DH=L^2,\qquad M^2H=RL^2.
\tag{1A-scale}
\]
The controlling vertices are
\[
 V_1=(5/18,1/3),\qquad V_2=(5/18,25/72),\qquad V_3=(7/24,1/3).
\]
For a prime $\ell$, put
\[
 \rho_\ell(n)=\1_{\ell\mid n}-\frac1\ell.
\]

A canonical row is $e=(r,k,m)$ with
\[
 r\asymp R\ \text{prime},\qquad1\le|k|\le K,\qquad m\asymp M,
 \qquad m'=m+kr\asymp M.
\]
Choose $w_0\pmod r$ with $mw_0\equiv-2\pmod r$.  Since $m'\equiv m\pmod r$, the same $w_0$ works for $m'$.
Define
\[
 \alpha=\frac{mw_0+2}{r},\qquad \beta=\frac{m'w_0+2}{r},
\]
\[
 A_e(t)=\alpha+mt,\qquad B_e(t)=\beta+m't.
\]

\begin{lemma}[Determinant-two row identity]\label{lem:det2}
For every row and integer $t$,
\[
 \boxed{m'A_e(t)-mB_e(t)=2k.}
\tag{1A-det}
\]
\end{lemma}
\begin{proof}
The $t$-terms cancel, and
\[
 m'\alpha-m\beta
 =\frac{m'(mw_0+2)-m(m'w_0+2)}r
 =\frac{2(m'-m)}r=2k.
\]
\end{proof}

Fix one common physical weight $W_D(t)=W(t/D)$.  For bounded prime coefficients $b_p,d_q$, define
\[
 P_e(t)=\sum_{p\asymp L}b_p\rho_p(A_e(t)),\qquad
 Q_e(t)=\sum_{q\asymp L}d_q\rho_q(B_e(t)),
\]
\[
 C_e(b,d)=\sum_{t\in\mathbb Z}W_D(t)P_e(t)Q_e(t).
\tag{1A-source}
\]
The common-weight condition is substantive: later row dependence must be generated by the transform, not chosen
independently on each row.

\section{Fourier normalization and the determinant shell}

The normalized pre-square coefficient has the schematic exact form
\[
 \widetilde C_{r,k,m}=
 \sum_{p\asymp L}b_p\sum_{\substack{q\asymp L\\q\ne p}}d_q
 \sum_h\omega_{m,p,q}(h)
 \e_{mq}\!\left(\frac{2hkp}{m+kr}\right),
\]
with
\[
 \omega_{m,p,q}(h)=\frac{H}{pq}\widehat W_D(h/pq)
 \e\!\left(-\frac{h\alpha}{pqm}\right).
\tag{1A-omega}
\]
The source norms satisfy
\[
 \sum_h|\omega_{m,p,q}(h)|^2\ll H,
 \qquad
 \sum_h|\omega_{m,p,q}(h)|\ll H.
\]
Additive reciprocity separates the phase, and the $h$-Poisson/Gauss completion produces the determinant shell
\[
 \boxed{ms-prn=2.}
\tag{1A-shell}
\]
If $\kappa_p(m)\pmod{pr}$ satisfies $m\kappa_p(m)\equiv2\pmod{pr}$, then all solutions are
\[
 s_j=\kappa_p(m)+jpr,
 \qquad
 n_j=\frac{m\kappa_p(m)-2}{pr}+jm.
\]
The quotient index has length $|j|\ll U=L/H$ and the corresponding amplitude has $\ell^2$ norm $\ll UX_1^{o(1)}$.

\section{The source-preserving Poisson--Bruhat input}

The next statement is a genuine Gate-1A hypothesis, not a consequence of the preceding algebra.

\begin{hypothesis}[One-sided source transport and absolute ceiling]\label{hyp:PB}
For every fixed $A_0$, the transported source splits into retained states and a discarded remainder whose Gate
energy is $O_{A_0}(X_1^{-A_0})$ times the natural scale.  Every retained state $\xi$ has row amplitude
\[
 \boxed{c_r(\xi)=B_\xi F_\xi(r/R),}
\tag{1A-PB}
\]
where $B_\xi$ is independent of $r$, $F_\xi$ is supported in a fixed compact subinterval of $(1,2)$, and for
every fixed $J$,
\[
 \sup_\xi\sup_t|\partial_t^JF_\xi(t)|\ll_JX_1^{o(1)}.
\]
Retained states satisfy a non-negligibility cutoff relative to the Fourier truncation.  The transform is one-sided
with $X_1^{o(1)}$ nuclear cost and supplies
\[
 T_{\rm abs}:=\sum_{r\asymp R\atop r\ \rm prime}
 \left(\sum_\xi|c_r(\xi)|\right)^2
 \ll M^2L^4X_1^{o(1)}.
\tag{1A-ceiling}
\]
Finally the physical and normalized sources satisfy
\[
 C_e=H^{-1}\widetilde C_e+R_e,
 \qquad
 \sum_e|R_e|^2\ll_{A_0}X_1^{-A_0}\frac{ML^4}{H}.
\tag{1A-bridge}
\]
\end{hypothesis}

The scale identity behind the transport budget is elementary:
\[
 \frac{MK(MH)^2}{RL^4}=1.
\tag{1A-budget}
\]

\section{Band-limited prime participation and the non-projective family energy}

Put $a_{r,\xi}=|c_r(\xi)|$ and
\[
 M_\xi=|B_\xi|\sup_{1\le t\le2}|F_\xi(t)|.
\]
Choose Fourier degree $N_0=R^{1/4}$.  Smoothness gives a trigonometric approximation
\[
 F_\xi(t)=P_\xi(t)+O_A(R^{-A}).
\]
Bernstein's inequality $\|P'\|_\infty\le N_0\|P\|_\infty$ yields an interval $I_\xi\subset[R,2R]$ of length
\[
 |I_\xi|\gg R^{3/4}
\]
on which $|c_r(\xi)|\gg M_\xi$.

\begin{hypothesis}[Fixed-state exclusion]\label{hyp:exclusion}
For each retained state, every row prime excluded from the generic non-projective branch divides one of
$X_1^{o(1)}$ fixed nonzero integers of polynomial height.
\end{hypothesis}

Huxley's prime theorem in intervals of length $x^{7/12+\epsilon}$ implies that each $I_\xi$ contains
$R^{3/4-o(1)}$ primes; deleting polynomially many fixed divisors costs only $X_1^{o(1)}$ primes.  Define
\[
 A_\xi=\max_{r\asymp R}a_{r,\xi},
 \qquad
 N_\xi=\frac{\sum_{r\asymp R}a_{r,\xi}}{A_\xi}.
\]
Then
\[
 A_\xi\asymp M_\xi,
 \qquad N_\xi\gg R^{3/4-o(1)}.
\tag{1A-part}
\]

Let $S$ be the number of prime rows, so $S=R^{1+o(1)}$.  Cauchy gives the finite participation inequality
\[
 \left(\sum_\xi A_\xi\right)^2\le \frac{S}{P^2}T_{\rm abs}
\tag{1A-Cauchy}
\]
whenever $N_\xi\ge P$ for all retained states.

Each state carries projective coordinates $(Z_\xi,L_\xi)$ and determinant
\[
 \Delta(\xi,\eta)=Z_\xi L_\eta-Z_\eta L_\xi.
\]
For $\Delta(\xi,\eta)\ne0$, a row prime can couple the pair only if it divides the fixed nonzero determinant,
so the pair occurs in only $X_1^{o(1)}$ rows.  With $P=R^{3/4-o(1)}$, one obtains
\[
 \boxed{E_{\rm np}\ll R^{-1/2+o(1)}T_{\rm abs}.}
\tag{1A-energy}
\]
This is the genuine BPP family-energy contraction.

\section{The outer square root: exact finite inequality and the source pin}

The finite four-cycle inequality is
\[
 \sum_r\|A_r\|_{\rm op}
 \le |\mathcal R|^{1/2}\left(\sum_r\operatorname{Tr}(A_r^*A_r)\right)^{1/2}.
\tag{1A-root-alg}
\]
To turn this into the physical Gate contribution one still needs the following source identification.

\begin{hypothesis}[Outer-root energy pin]\label{hyp:rootpin}
Let $G_{\rm np}$ be the normalized Gate contribution of the non-projective transported branch.  Then
\[
 \boxed{G_{\rm np}\ll M^2L^4X_1^{o(1)}
 \left(\frac{E_{\rm np}}{T_{\rm abs}}\right)^{1/2}.}
\tag{1A-root-pin}
\]
\end{hypothesis}

Combining with (1A-energy) gives
\[
 G_{\rm np}\ll R^{-1/4+o(1)}M^2L^4.
\]
The target comparison is
\[
 MR^{-1/4}\le H
 \iff \frac54a+2b-1\ge0.
\]
At $V_1,V_2,V_3$ the margins are respectively
\[
 \boxed{\frac1{72},\qquad\frac1{24},\qquad\frac1{32}.}
\tag{1A-margins}
\]
The binding margin is $1/72$: any later fixed-power loss must be explicitly budgeted against it.

\section{The projective branch and the full state diagonal}

If $\Delta(\xi,\eta)=0$, the exact collision relation is
\[
 Z_\xi L_\eta=Z_\eta L_\xi.
\]
For $L\ne0$, the correct equivalence class is the ratio $Z/L$.  For vectors $v_{Z,L}=A_Z\otimes B_L$, define
\[
 C_\rho=\sum_{(Z,L):Z/L=\rho}v_{Z,L}.
\]
Then the finite identity is
\[
 \boxed{
 \sum_\rho\|C_\rho\|^2
 =\sum_{(Z_1,L_1),(Z_2,L_2)}
 \1_{Z_1L_2=Z_2L_1}
 \operatorname{Re}\langle v_{Z_1,L_1},v_{Z_2,L_2}\rangle.}
\tag{1A-proj-id}
\]
Since $\Delta(\xi,\xi)=0$, the entire state diagonal is in this branch.

\begin{hypothesis}[Projective/diagonal source routing]\label{hyp:proj}
The physical projective source is routed into positive primitive ratio classes of multiplicity $X_1^{o(1)}$;
axis and non-primitive/rank-loss states are routed to explicit exceptional blocks.  Moreover
\[
 \boxed{\sum_{(Z,L)}\|A_Z\|^2\|B_L\|^2\ll MHL^4X_1^{o(1)},}
\tag{1A-diag}
\]
and every exceptional block is bounded at the same target scale.
\end{hypothesis}

This is a load-bearing conditional statement: it controls a branch containing the complete diagonal at the
final target scale.  It is not generated by the off-diagonal BPP contraction.

\section{Recombination and the conditional Gate-1A theorem}

\begin{hypothesis}[Literal source recombination estimate]\label{hyp:rec}
The quotient-cell/source recombination error satisfies
\[
 \boxed{E_{\rm rec}\ll U^{-2+o(1)}M^2L^4.}
\tag{1A-rec}
\]
\end{hypothesis}
Since $U=L/H$ and $DH=L^2$,
\[
 \frac{U^{-2}M^2L^4}{MHL^4}=\frac{M}{D},
\]
and the margins at $V_1,V_2,V_3$ are $1/18,1/18,1/24$.

\begin{theorem}[Conditional canonical Gate-1A estimate]\label{thm:G1A}
Assume Hypotheses~\ref{hyp:PB}, \ref{hyp:exclusion}, \ref{hyp:rootpin}, \ref{hyp:proj}, and \ref{hyp:rec}.
Then throughout the polytope generated by $V_1,V_2,V_3$,
\[
 \boxed{\sum_e|C_e(b,d)|^2\ll \frac{ML^4}{H}X_1^{o(1)}.}
\tag{1A-main}
\]
Equivalently, in normalized form,
\[
 \sum_e|\widetilde C_e(b,d)|^2\ll MHL^4X_1^{o(1)}.
\]
\end{theorem}
\begin{proof}
The non-projective branch follows from (1A-energy) and the outer-root pin.  The complete projective branch,
including the diagonal, follows from the exact ratio-class identity and Hypothesis~\ref{hyp:proj}.  The source
recombination error is acceptable by Hypothesis~\ref{hyp:rec} and its exponent margin.  The physical/normalized
bridge in Hypothesis~\ref{hyp:PB} gives the displayed physical estimate.
\end{proof}

\section{What Gate 1A does and does not give}

The theorem above is an analytic theorem for the canonical common-weight operator.  It does not prove that every
packet of an arbitrary global sieve decomposition is of this form.  A source-exhaustive compiler would need a
decomposition
\[
 S_{\rm actual}=\sum_{\lambda\in\Lambda_{\rm gen}}a_\lambda D_\lambda U_\lambda S^\sharp_\lambda+S_{\rm exc},
 \qquad
 \sum_\lambda|a_\lambda|\|D_\lambda\|_{\rm op}\le X_1^{o(1)},
\tag{1A-global}
\]
with exact packet identity, multiplicity control and exceptional routing.  This is a separate theorem.

If Gate 1A is to be used for 287, a further scale dictionary is mandatory because Gate 1A uses
$m\asymp X_1^{1/3}$ while the Ford--Maynard/287 Type-II variable lies in
$(X/2)^\varepsilon<m\le X^{1/6-\varepsilon}$.  One must specify $X_1$ as a function of the 287 scale, identify
the physical bilinear variable, recheck $R,L,H,K,U,D$, and preserve the $1/72$ margin.  No identity follows
merely from the common letter $X$.

\part{Gate 1B: exact conductor reduction, late parent-operator reorganisation, and the pair-modulus frontier}


\section{Gate-1B front layer: fixed-shift determinant-2 algebra}
\label{sec:g1b-det2-front}

Before the H7 conductor calculation there is a source-native determinant layer which is independently checkable.
It is useful to write it explicitly because it identifies which cancellation resources really exist.

\subsection{M\"obius sign collapse at a distinguished prime}

Suppose $q=dp$, $p$ is prime and $(d,p)=1$.  Multiplicativity of the M\"obius function gives
\[
 \mu(q)=\mu(d)\mu(p)=-\mu(d),
\]
so
\[
 \boxed{\mu(d)=-\mu(q).}
\tag{1B-F1}
\]
On squarefree support, if $p\mid q$ and $d=q/p$, then $(d,p)=1$ automatically.  Hence all admissible distinguished
prime representations of a fixed squarefree $q$ carry the \emph{same} M\"obius sign $-\mu(q)$.  There is no
same-$q$ cancellation obtained merely by summing over the distinguished prime $p$.

For a finite weighted cell
\[
 \lambda_{D,P}(q)=\sum_{\substack{dp=q\\p\in P}}\mu(d)\log p,
 \qquad
 L_{D,P}(q)=\sum_{\substack{dp=q\\p\in P}}\log p,
\]
this gives the exact squarefree identity
\[
 \boxed{\lambda_{D,P}(q)=-\mu(q)L_{D,P}(q).}
\tag{1B-F2}
\]

\subsection{Complementary divisor and determinant-2 normal form}

For positive integers $d,p,u,v$,
\[
 dp\mid uv+2
 \quad\Longleftrightarrow\quad
 \exists\ell\ge1:\ uv+2=dp\ell.
\tag{1B-F3}
\]
Writing $q=dp$, the incidence is equivalently
\[
 \boxed{q\ell-uv=2.}
\tag{1B-F4}
\]
The complementary divisor is unique and equals $(uv+2)/q$.

Any common divisor of $u$ and $\ell$ divides both $uv$ and $q\ell$, hence divides their difference $2$:
\[
 \boxed{(u,\ell)\mid2.}
\tag{1B-F5}
\]
Similarly
\[
 \boxed{(v,q)\mid2.}
\tag{1B-F6}
\]
so on the odd sector these pairs are coprime.  This is the precise reason ordinary modular inverses become
available there.

If $(q_1,v_1)$ and $(q_2,v_2)$ lie on the same fixed $(u,\ell)$ determinant-2 line,
\[
 q_i\ell-uv_i=2,
\]
then subtraction yields
\[
 \boxed{u(v_2-v_1)=\ell(q_2-q_1).}
\tag{1B-F7}
\]
Thus the solution set is an affine line; translation along the line preserves the determinant equation.

\subsection{Source-native HFMV determinant identity}

For two incidences
\[
 uv_i+2=d_ip_i\ell_i\qquad(i=1,2),
\tag{1B-F8}
\]
one has the exact determinant identity
\[
 \boxed{
 v_2d_1p_1\ell_1-v_1d_2p_2\ell_2=2(v_2-v_1).}
\tag{1B-F9}
\]
Proof: substitute $d_ip_i\ell_i=uv_i+2$; the $uv_1v_2$ terms cancel.
If $v_1=v_2$, then
\[
 d_1p_1\ell_1=d_2p_2\ell_2.
\tag{1B-F10}
\]
Conversely, when $u\ne0$, equality of these products forces $v_1=v_2$ by (1B-F8).

A finite B1 multiplicity statement also follows in the separated-prime regime: when $p_1\ne p_2$ are primes,
$|h_1|,|h_1'|\le H$ and $2H<p_1$, the determinant congruence has at most one admissible $h_1$ in its interval.
This supplies multiplicity one for that finite fibre.  What it does \emph{not} supply is the generic analytic
off-diagonal estimate.

The source-native HFMV reduction therefore isolates the following analytic obligations:
\[
 \boxed{\begin{array}{c}
 \text{generic off-diagonal determinant variance}\\
 +\ \text{diagonal divisor bound}\\
 +\ \text{small proper-gcd sectors}\\
 +\ \text{literal source centering}
 \end{array}}
\tag{1B-F11}
\]
The exact determinant algebra above is unconditional; the four displayed estimates are separate analytic/source
theorems.

\subsection{Near-top exponent geometry}

In the high-modulus window $13/18\le\omega\le8/9$, write
\[
 Q_e=\omega,\qquad R_e=1-\omega,
\qquad H_e=Q_e-R_e=2\omega-1.
\tag{1B-F12}
\]
Then
\[
 \frac19\le R_e\le\frac5{18},\qquad 2R_e<Q_e,
\tag{1B-F13}
\]
and at the binding endpoint $\omega=13/18$,
\[
 H_e=\frac49,
 \qquad
 \frac{H_e}{Q_e}=\frac8{13}.
\tag{1B-F14}
\]
The associated near-top Kloosterman exponent ledger records
\[
 \frac13\frac8{13}-\frac15=\frac1{195}>0.
\tag{1B-F15}
\]
Again, (1B-F15) is an exponent margin, not by itself a Kloosterman theorem.

A useful exact local character identity is, for a nonzero field element $u$,
\[
 uv+2=u\left(v+2u^{-1}\right),
\tag{1B-F16}
\]
hence for every multiplicative character $\chi$,
\[
 \boxed{\chi(uv+2)=\chi(u)\chi(v+2u^{-1}).}
\tag{1B-F17}
\]
For odd prime modulus the map $u\mapsto2u^{-1}$ is injective on nonzero residues.  This is a finite
reparametrization tool; it does not provide cancellation by itself.


\section{High-order decomposition and the conductor split}

The high-order source is expanded schematically as
\[
 a_{ubv}=\prod_{i=1}^9(f_i+\delta_i)
 =\sum_{J\subset[9]}\delta_Jf_{J^c}.
\tag{1B-exp}
\]
An order-$j$ term carries a determinant shell of the form
\[
 C_Ja_J-q\ell=-2,
 \qquad C_J\asymp Y^j,
 \qquad a_J\asymp Y^{9-j}.
\tag{1B-shell}
\]
The exact centering, CRT, Gauss and determinant identities which feed the lower orders have been written out in
Sections~\ref{sec:g01-front} and~\ref{sec:g1b-det2-front}.  Their role is to route the lower-complexity pieces
and expose the fixed-shift determinant geometry; any analytic estimate beyond those identities must still be
stated separately.  Order $6$ is regrouped with the order-$5$ geometry.  H7 and H8 admit one-dimensional
character/Fourier representations, while H9 is the pure-defect child.

For H7, write the induced character modulus
\[
 q=ce,
 \qquad c\asymp C,
 \qquad q\asymp Q.
\]
The corrected coefficient-state count is
\[
 \sum_{c,\chi^*,e}|\beta(ce)|^2\ll QCX^{o(1)}.
\tag{1B-state}
\]
The coefficient-blind estimate has relative size
\[
 \frac{|H_7(C)|}{Y^9}\ll
 \begin{cases}
 C/Q,&C\le Y^4,\\[1mm]
 C^2/(QY^4),&C\ge Y^4.
 \end{cases}
\tag{1B-blind}
\]
The formulae cross at $C=Y^4$, while the large-$C$ estimate loses all saving at
\[
 \boxed{C_0=Q^{1/2}Y^2.}
\tag{1B-C0}
\]
Thus $C_0$, not $Y^4$, is the high-conductor zero-saving threshold.

\section{Exact high-conductor source split and spent M\"obius sign}

On the squarefree source,
\[
 \beta(q)=\sum_{dp=q}\mu(d)\log p,
\]
with physical restrictions $p>V$ and $q/p>U$.  In high conductor
\[
 e=q/c\le Y^2,
 \qquad V=Y^{5/2-o(1)},
\]
so every distinguished source prime satisfies
\[
 p>V>e.
\]
Hence the lane $p\mid e$ is empty.  Write
\[
 c=pc_0,
 \qquad d=c_0e,
\]
and on the clean squarefree sector
\[
 \boxed{\beta(ce)=\mu(e)\rho_{C,e}(c),
 \qquad
 \rho_{C,e}(c)=\sum_{pc_0=c}\mu(c_0)\log p,}
\tag{1B-split}
\]
with the physical cell restrictions retained.

Let $\chi^*$ be primitive modulo $c$ and let $\chi$ be its induction to $ce$.  For squarefree $(c,e)=1$,
\[
 \tau_{ce}(\chi)=\mu(e)\chi^*(e)\tau_c(\chi^*).
\tag{1B-Gauss}
\]
Multiplying by (1B-split),
\[
 \boxed{\beta(ce)\tau_{ce}(\chi)
 =\rho_{C,e}(c)\chi^*(e)\tau_c(\chi^*).}
\tag{1B-spent}
\]
Thus the short-cofactor sign $\mu(e)$ is algebraically spent and cannot be counted again as an independent
cancellation resource.

\section{Prime-character collapse}

Write $c=pc_0$ with $(p,c_0)=1$ and factor a primitive character as $\chi=\chi_p\chi_0$.  Put
\[
 N=r_1\cdots r_7x\asymp Y^8.
\]
The $p$-local argument has the form
\[
 A_p\equiv c_0eN(2h)^{-1}\pmod p.
\]

\begin{lemma}[Prime-character collapse]\label{lem:charcollapse}
For $A\in\mathbb F_p^\times$,
\[
 \boxed{
 \sum_{\chi_p\ne\chi_{0,p}}\frac{\tau_p(\chi_p)}{\sqrt p}\chi_p(A)
 =\frac{p-1}{\sqrt p}\e_p(\bar A)+\frac1{\sqrt p}.}
\tag{1B-char}
\]
\end{lemma}
\begin{proof}
Summing over all multiplicative characters gives $(p-1)\e_p(\bar A)$.  The principal character has Gauss sum
$-1$, so removing it adds $1$.  Divide by $\sqrt p$.
\end{proof}

The main term converts the reciprocal character dependence on $h$ into the linear additive phase
\[
 \e_p(2hc_0eN).
\tag{1B-linear}
\]
The correction is smaller than the main local multiplier by $p^{-1}$ and is power-negligible because
$p>V=Y^{5/2-o(1)}$, provided the stated source ceiling is available.

\section{Hybrid Poisson and the primitive projector}

The long Fourier slot has length
\[
 H=Q/Y,
 \qquad \frac cH=\frac Ye.
\]
Hence the dual $h$-frequency window has length $\asymp Y/e<p$.

\begin{lemma}[Hybrid Poisson]\label{lem:hybrid}
Let $p\nmid c_0$ and let $\chi_0$ be primitive modulo $c_0$.  With the Fourier convention
$\widehat v(\xi)=\int v(x)\e(-x\xi)\,dx$,
\[
 \boxed{
 \sum_{h\in\mathbb Z}v(h/H)\chi_0(h)\e_p(ah)
 =\frac H{c_0}\tau_{c_0}(\chi_0)
 \sum_{\substack{m\in\mathbb Z\\m\equiv ac_0\ (p)}}
 \chi_0(-m\bar p)\widehat v\!\left(\frac{mH}{pc_0}\right).}
\tag{1B-Poisson}
\]
\end{lemma}
\begin{proof}
Apply Poisson summation in residue classes modulo $pc_0$.  CRT separates the finite Fourier transform.  The
$p$-factor enforces $m\equiv ac_0\pmod p$, while the $c_0$-factor is the twisted Gauss sum; the factor $p$ from
the $p$-sum cancels the $p$ in the Poisson denominator.
\end{proof}

For $a=2c_0eN\pmod p$, the surviving dual frequency satisfies
\[
 \boxed{meN\equiv2\pmod p,
 \qquad |m|\ll (Y/e)X^{o(1)}<p.}
\tag{1B-dualfreq}
\]

\begin{lemma}[Primitive-character projector]\label{lem:projector}
For odd squarefree $c_0$ and $(A,c_0)=1$,
\[
 \boxed{\sum_{\chi_0\ (\mathrm{mod}\ c_0)}^*\chi_0(A)
 =\sum_{\substack{d\mid c_0\\d\mid A-1}}\mu(c_0/d)\varphi(d).}
\tag{1B-proj}
\]
Consequently
\[
 \boxed{\frac{\mu(c_0)}{c_0}\sum_{\chi_0}^*\chi_0(A)
 =\frac1{c_0}\sum_{\substack{d\mid c_0\\d\mid A-1}}\mu(d)\varphi(d).}
\tag{1B-proj2}
\]
\end{lemma}
\begin{proof}
This is M\"obius inversion over the conductor.  For squarefree $c_0$,
$\mu(c_0)\mu(c_0/d)=\mu(d)$.
\end{proof}

With $A=meN/2$, the prime congruence and primitive projector combine into
\[
 \boxed{pd\mid meN-2.}
\tag{1B-detproj}
\]

\section{The H7 dual determinant and fixed-power repair}

After the exact reductions, the H7 main term is
\[
 \begin{aligned}
 D_7(C)=&\sum_{pc_0\asymp C\atop p>V}
 \sum_{e\asymp Q/C}
 \log p\left(1-\frac1p\right)
 \sum_{N\asymp Y^8}a_7(N)
 \sum_{|m|\ll(Y/e)X^{o(1)}\atop p\mid meN-2}
 \widehat v(me/Y)\\
 &\hspace{3cm}\times
 \frac1{c_0}
 \sum_{d\mid c_0\atop d\mid meN-2}\mu(d)\varphi(d).
 \end{aligned}
\tag{1B-D7}
\]
Here $a_7(N)$ is the literal seven-defect plus one-model source coefficient.

\begin{hypothesis}[H7 source norm and cell multiplicity]\label{hyp:H7norm}
The physical source satisfies
\[
 \sum_{N\asymp Y^8}|a_7(N)|\ll Y^8X^{o(1)},
\]
$\widehat v$ is rapidly decreasing with bounded supremum, the omitted dyadic/coprimality selectors have total
multiplicity $X^{o(1)}$, and the primitive-character family above is the literal family generated by the source.
\end{hypothesis}

\begin{proposition}[H7 fixed-power repair]\label{prop:H7power}
Under Hypothesis~\ref{hyp:H7norm},
\[
 \boxed{|D_7(C)|\ll Y^9X^{o(1)}.}
\tag{1B-H7power}
\]
\end{proposition}
\begin{proof}
For fixed $(e,m,N)$, the nonzero integer $meN-2$ has only $O(1)$ prime divisors $p>V=Y^{5/2-o(1)}$.
After $p$ is fixed, write $c_0=ds$.  The absolute divisor projector is bounded by
\[
 \sum_{d\mid meN-2}\sum_{s\asymp C/(pd)}\frac{\varphi(d)}{ds}\ll X^{o(1)},
\]
because the $s$-sum is harmonic on a dyadic interval and the divisor count is $X^{o(1)}$.  Summing
$e\asymp E$ and $|m|\ll Y/e$ costs $O(YX^{o(1)})$; use the $\ell^1$ norm of $a_7$.
\end{proof}

Thus the previous coefficient-blind fixed-power loss $Q/Y^4$ has been recovered.  The target, however, is
\[
 \boxed{D_7(C)\ll_A Y^9(\log X)^{-A}.}
\tag{H7-open}
\]
This arbitrary-log saving is still open.  The full divisor term $d=c_0$ reconstructs the original two-model
determinant skeleton, so the character/Poisson transform is self-dual and cannot simply be iterated to closure.


\section{Historical H7 natural-scale status and the late short-short promotion}
\label{sec:g1b-h7-promotion}

The natural-scale estimate (\ref{prop:H7power}) remains an important reduction: it removed the former
coefficient-blind fixed-power deficit.  It is no longer the controlling final status for the entire H7 analysis.
In the \emph{short-short} exponent region
\[
 \boxed{\alpha<\frac49,\qquad \beta<\frac49,}
\tag{1B-SS}
\]
the late Gate-1B audit isolates a joint-prime character packet.  The deterministic compiler has the following
conditional form.

\begin{hypothesis}[H7 short-short common-sequence and energy inputs]\label{hyp:H7SS}
After the source-exact H7 character reduction, the short-short packet is represented on one common sequence;
its coefficient energy is bounded at the audited source scale; and the required multiplicative large-sieve
estimate applies uniformly to that common sequence with the physical dyadic/coprimality restrictions retained.
\end{hypothesis}

\begin{proposition}[H7 short-short joint-prime capacity]\label{prop:H7SS}
Under Hypothesis~\ref{hyp:H7SS}, the H7 short-short contribution gains the relative fixed power
\[
 \boxed{Y^{-1/2}=X^{-1/18}}
\tag{1B-H7SS}
\]
against the natural target scale.
\end{proposition}

The statement above is a \emph{conditional analytic compiler}: the capacity is positive, but the source pins and common-sequence pins remain part of the theorem.  It should therefore not be described as an unconditional H7 closure.
The high-prime complement is disjoint from (1B-SS) and is not absorbed by this compiler.

The corresponding H8 audit reports the same fixed-power capacity after the analogous source-energy calculation.
For publication purposes the status is deliberately asymmetric:
\[
 \boxed{\begin{array}{ll}
 \text{H7 short-short:}&\text{conditional compiler with }X^{-1/18}\text{ capacity},\\
 \text{H8 short-short:}&\text{same research capacity, but source/robustness interface still open}.
 \end{array}}
\tag{1B-H78-status}
\]
Thus the earlier H7 arbitrary-log determinant (H7-open) remains a valid historical source-specific target, but it
is no longer the sole controlling description of the late H7 frontier.

\section{The full $4/9$ complement and the square-root transition}
\label{sec:g1b-49-complement}

The complement of the short-short region is organised by the exponents of the large source primes.  At programme
level one separates a high-$d$ wing and a high-$p$ wing.  Joint-prime machinery has positive capacity on the
high-$d$ side and on the high-$p$ side away from the transition where the large prime approaches $X^{1/2}$.
The square-root strip is not a formal consequence of the short-short compiler and must be routed separately.

The purpose of this decomposition is logical as well as analytic: a bound in one exponent cell cannot be moved
into a disjoint cell merely because both cells arise from the same order-$j$ source.  The square-root transition is
therefore placed inside the shifted-quotient parent below.

\section{The Full-Nine shifted-quotient parent}
\label{sec:g1b-fullnine}

The high-prime switch and the pure H9 specialization can be embedded in one parent operator.  Its source-exact
schematic form is
\[
 \boxed{
 \mathcal S_{\rm SQ}
 =\sum_{mp\asymp X}\gamma(m)\lambda(p)\log p\;c_{\rm res}(mp-2).}
\tag{1B-FullNine}
\]
Here $\gamma$ and $\lambda$ retain the actual source-generated convolution, M\"obius/sign and support structure;
they are not replaced by arbitrary independent coefficients.  The high-prime $M$-switch is one subrange of
(1B-FullNine), while H9 is the pure-nine-defect specialization.  Type-I or joint-large-sieve methods have capacity
on strongly unbalanced sides.  The balanced square-root strip is the difficult transition.

Equation (1B-FullNine) is a \emph{parent operator}, not a theorem asserted closed.  A successful estimate must
control the complete source-generated coefficient family on the relevant subranges rather than a single fixed
coefficient choice.

\section{The QK56 full-covariance parent and its internal promotion}
\label{sec:g1b-qk56}

A second parent retains the complete covariance of the one-copy QK output.  If $Q(q;h,k)$ denotes the exact
one-copy output after the local transforms, the full covariance has schematic form
\[
 \boxed{
 \mathcal C_{\rm QK}
 =\sum_{q_1,q_2}\sum_{h_1,k_1,h_2,k_2}
 A(q_1,q_2;h_1,k_1,h_2,k_2)
 Q(q_1;h_1,k_1)Q(q_2;h_2,k_2).}
\tag{1B-QK56}
\]
The decomposition
\[
 \boxed{q_1=q_2\quad\text{versus}\quad q_1\ne q_2}
\tag{1B-QKsplit}
\]
is the same-modulus diagonal versus cross-modulus off-diagonal of this one covariance parent.

\subsection{Same-modulus covariance: character diagonal and cross-character Gram}

For $q_1=q_2=q$, the source audit expands the remaining covariance into
\[
 \boxed{
 \sum_{q\asymp Q}|\beta(q)|^2
 \sum_{\chi_1,\chi_2\ (\mathrm{mod}\ q)}
 \widehat R_q(\chi_1)\widehat R_q(\chi_2)W_q(\chi_1,\chi_2)
 =G_{\rm diag}+G_{\rm off}.}
\tag{1B-QK-same}
\]
The character diagonal $\chi_1=\chi_2$ carries the large $HK/q$ contribution and is routed to the existing
coefficient-blind residue-energy/JQ7 estimate.  The cross-character term is a centred modular-hyperbola
covariance.  The current internal audit estimates it through the corresponding square-root Weil/Kloosterman
discrepancy and records the worst signed programme margin
\[
 \boxed{X^{-5/36+o(1)}.}
\tag{1B-QK-margin}
\]
Accordingly the programme-level status is
\[
 \boxed{\mathrm{SAMEQ\!\mbox{-}QK56\!\mbox{-}CROSSCHAR\!\mbox{-}GRAM45}:
 \text{ internal analytic pass,}}
\tag{1B-QK-same-status}
\]
subject to the literal QK56/S1/S2 source transcription and a final hostile audit.  This is an internal analytic
claim, not a published theorem and not a Lean-closed analytic leaf.

\subsection{Cross-modulus no-wrap branch: source-spread certificate}

For $q_1\ne q_2$, the late source audit uses the no-wrap regime: after the physical source parameter $\Theta$
is fixed, the residues $B_1\pmod{q_1}$ and $B_2\pmod{q_2}$ are fixed.  The remaining product multiplicity is
then measured by the push-forward of the source multiplier.  The internal audit records that the relevant
$\ell^2/\ell^1$ spread lies within the required source budget.  We record
\[
 \boxed{\mathrm{CROSSQ\!\mbox{-}THETA\!\mbox{-}SPREAD45}:
 \text{ internal finite/source pass,}}
\tag{1B-QK-cross-status}
\]
again subject to final literal transcription of the physical source.

Combining the two children gives the programme status
\[
 \boxed{\mathrm{QK56\!\mbox{-}FULL\!\mbox{-}COVARIANCE45}:
 \text{ analytic pass modulo literal source pins}.}
\tag{1B-QK-full-status}
\]
This is a substantial promotion relative to Draft V9: QK56 is no longer the controlling major analytic open,
provided the internal same-$q$/cross-$q$ audits survive source transcription and hostile checking.

\section{Canonical Gate-1B zero mode and the historical comparison fork}
\label{sec:g1b-zero}

The canonical additive zero mode is not part of the analytic frontier.  The finite/algebraic audit gives three
exact statements:
\begin{enumerate}[label=(\roman*)]
\item the canonical additive-character baseline is recovered from the standard character system;
\item the canonical discrepancy has zero additive Fourier frequency;
\item if a historical formulation retains an auxiliary comparison $E(q)$, then
\[
 \boxed{S_{\rm hist}=S_{\rm can}-R_E.}
\tag{1B-zero-hist}
\]
\end{enumerate}
A perturbation of $E(q)$ changes $R_E$ while leaving the canonical nonzero-frequency packet unchanged.  Hence
\[
 \boxed{\text{canonical zero mode: algebraic PASS};\qquad
        \text{historical }E(q)\text{ identification: separate source fork}.}
\tag{1B-zero-status}
\]

\section{The balanced shifted child: recover the source before analysis}
\label{sec:g1b-shifted-child}

The balanced shifted-quotient branch must not be modelled as four independent arbitrary multipliers.  Its
source action is cyclic.  Write the physical packet variables as $B_1,\ldots,B_4$ and let $\Theta$ denote the
literal source map.  The multiplier relation has the form
\[
 \boxed{a_j=\Theta(B_j,B_{j+1})\qquad(j\bmod 4),}
\tag{1B-shift-cycle}
\]
and the associated character moment has schematic source form
\[
 \boxed{
 \sum_{B_1,\ldots,B_4}W(B)
 \sum_h W_B(h)\,
 \chi\!\left(
 D\bigl(\Theta(B_1,B_2),\Theta(B_2,B_3),
        \Theta(B_3,B_4),\Theta(B_4,B_1);h\bigr)
 \right).}
\tag{1B-shift-moment}
\]
The four-cycle discriminant algebra is already available, but the literal pre-support-enlargement $TT^*$ source
formula has not yet been recovered.  Therefore the first required action is not another black-box estimate; it is
\[
 \boxed{\mathrm{SHIFT\!\mbox{-}TTSTAR\!\mbox{-}LITERAL\!\mbox{-}SOURCE\!\mbox{-}PIN}.}
\tag{1B-shift-source-pin}
\]
This source theorem must print, without enlargement or abstraction:
\begin{enumerate}[label=(\alph*)]
\item the literal ranges of $B_1,\ldots,B_4$;
\item the exact cyclic map $\Theta$ and all repeated-prime/collision constraints;
\item the $h$-range, smooth weight, prefactor and normalization;
\item the physical target before and after normalization;
\item all zero/nonunit/shared-prime/degenerate strata and their routing;
\item the multiplicity of every source packet under the $TT^*$ expansion.
\end{enumerate}
Only after this literal source recovery is the analytic child well posed:
\[
 \boxed{\mathrm{SHIFT\!\mbox{-}SOURCE\!\mbox{-}LINKED\!\mbox{-}CHAR45}.}
\tag{1B-shift-open}
\]
Conditional on the QK56 promotion surviving final audit, (1B-shift-open) is the current first major analytic
Gate-1B candidate.

\section{Technology cards for the shifted source-linked child}
\label{sec:g1b-shift-tech}

The following published technologies are relevant only after the literal source has been reduced to their stated
geometries.
\begin{enumerate}
\item Fouvry--Shparlinski estimate weighted multiplicative-character sums of the $2\times2$ determinant
$ad-bc$ over four interval variables, with power saving once the interval length is at least $p^{1/8+\varepsilon}$.
This is a serious candidate if the cyclic discriminant reduces source-exactly to determinant geometry.
\item Shparlinski's multivariate character/exponential-sum results allow monomial and reciprocal-monomial
kernels with bounded one-variable weights.  They become relevant if the $\Theta$-dictionary produces such a
kernel.
\item Bourgain--Garaev reciprocal sumset and multilinear Kloosterman technology is a secondary candidate after
an actual reciprocal-product dictionary has been manufactured.
\item General trace-function bilinear theorems require nontrivial geometric-monodromy hypotheses.  A rank-one
Kummer/Legendre character is not automatically covered merely because it is a trace function, so such a theorem
is not used here as a black box.
\item Blomer--Pascadi remains a fixed-Kloosterman backend-capacity card and likewise does not by itself estimate
the source-linked cyclic family.
\end{enumerate}
These are technology cards, not proof substitutions.

\section{Historical pair-modulus umbrella and the updated controlling frontier}
\label{sec:g1b-pairmod}

The earlier pair-modulus formulation remains useful as an umbrella.  Both the balanced shifted $TT^*$ child and
the QK cross-modulus child lead to product moduli, and fixed-multiplier Kloosterman technology has genuine
positive capacity.  However the Draft-10 source split promotes the QK56 covariance internally and leaves the
shifted child as the principal unresolved analytic family.  Consequently the historical labels
\[
 \mathrm{SOURCE\!\mbox{-}QUADRATIC\!\mbox{-}MULTIPLIER45},\qquad
 \mathrm{FM\!\mbox{-}PERRON\!\mbox{-}PAIRMOD\!\mbox{-}SOURCE\!\mbox{-}MULT45}
\]
are retained as umbrella/source-grammar labels only; they are no longer the controlling first-open theorem.

\section{Updated Gate-1B closure ledger}
\label{sec:g1b-late-ledger}

The current Gate-1B status is therefore:
\begin{center}
\begin{tabular}{p{0.38\textwidth}p{0.52\textwidth}}
\toprule
Component & Current status\\
\midrule
Orders 1--4 / low-order geometry & structural material available\\
H6 regroup geometry & structural material available\\
H7 short-short & conditional joint-prime compiler; $X^{-1/18}$ capacity\\
H8 short-short & same research capacity; robustness/source pin open\\
high-$d$/high-$p$ complement & routing/capacity programme; square-root strip routed to shift parent\\
QK56 same-$q$ cross-character Gram & internal analytic pass; source pins pending\\
QK56 cross-$q$ source spread & internal finite/source pass; transcription pending\\
QK56 full covariance & analytic pass modulo literal source pins\\
canonical additive zero mode & exact algebraic pass\\
historical $E(q)$ residual & source-identification fork if retained\\
SHIFT-$TT^*$ literal source pin & \textbf{SOURCE FIRST ACTION -- OPEN}\\
SHIFT source-linked character child & \textbf{FIRST MAJOR ANALYTIC OPEN}\\
S1/S2 literal source pins & narrow source interfaces remain\\
degeneracy/collision/packet census & finite routing/reassembly remains\\
\textbf{Gate 1B} & \textbf{NOT CLOSED}\\
\bottomrule
\end{tabular}
\end{center}
Conditional on the QK56 internal audit surviving final source transcription, the current Gate-1B picture has one
principal analytic child together with literal source and finite reassembly work.

\part{Source-exhaustive Full Type II and Gate 2}


\section{The F1 / F2 / F3 assembly that precedes Full Type II}
\label{sec:F123}

The phrase ``Full Type II'' hides three distinct families in the programme.  Their current dependency structure
can be written without any formal-system terminology.

\subsection{F1: global centering and the routable long-M\"obius sector}

Suppose the physical sequence is decomposed into finitely many pieces
\[
 a(n)=\sum_P A_P(n),
 \qquad
 A_P(n)=MT_P(n)+OD_P(n),
\]
and define
\[
 b(n):=\sum_PMT_P(n),\qquad w(n):=a(n)-b(n).
\]
Then identically
\[
 \boxed{w(n)=\sum_POD_P(n).}
\tag{F1-center}
\]
This is the exact F1 global-centering identity.  It does not prove that the chosen $MT_P$ are the correct physical
main terms; that is the comparison/source problem.

A long-M\"obius F1 piece is declared \emph{routable} when it contains a genuine M\"obius block of
mesoscopic-or-longer length and a smooth $(1)$-block of exponent $w$ satisfying
\[
 w\ge w_*(\mu)+\delta,
 \qquad \delta>0,
\tag{F1-route}
\]
with M\"obius cancellation used exactly once before an absolute value and with no nonprincipal character placed
on that M\"obius block.  The routing inequality is equivalent to the widened wedge
\[
 \boxed{122\mu+162\theta_j<1,
 \qquad \theta_j=\frac12-w.}
\tag{F1-wedge}
\]
Given a routed-F3 theorem on that wedge, such an F1 piece inherits an arbitrary logarithmic saving.  The missing
global step is to show that \emph{every} long-M\"obius F1 piece either satisfies this route or is controlled by a
separate theorem.  Hence routable F1 migration is not yet full F1.

\subsection{F2: prime/unit sector, principal sectors and reciprocal tensor}

The F2 family splits first into a prime/unit sector and excluded strata:
\[
 \boxed{\begin{array}{c}
 \text{principal characters},\ \text{nonunits},\ \text{equal primes},\\
 \text{prime powers},\ \text{repeated primes},\ \text{gcd strata}
 \end{array}}
\tag{F2-strata}
\]
On the prime/unit sector the required transformation is a double-Mellin / multiplicative-character expansion.
A complete source-exact formula for that transformation is itself still an open source theorem.  Thus the
statement to be proved is not replaced by a reference: one must derive the prime/unit character expansion from
the physical F2 source, including principal-character subtraction, unit restrictions, repeated/equal-prime
strata and all normalization factors, before the sector decomposition below is valid.

After the prime/unit reduction, the character sectors are
\[
 PP,\\ PN,\\ NP,\\ NN.
\]
The required F2 dependency is
\[
 \boxed{\begin{array}{c}
 \text{prime/unit reduction}\\
 +\ PP\text{ main}+PN\text{ single-spectral}+NP\text{ single-spectral}\\
 +\ NN\text{ reciprocal-character tensor}\\
 +\ \text{composite/gcd reassembly}\\
 \Downarrow\\
 \text{Full F2}
 \end{array}}
\tag{F2-full}
\]
The $NN$ reciprocal-character tensor estimate is a genuine open analytic input.  A structured tensor theorem is
Ford-ready only if it simultaneously preserves the fixed shift, arbitrary divisor-bounded coefficients, rank-one
structure and composite moduli.  A theorem missing any of those four properties does not fill the F2 socket.

\subsection{F3: fixed-depth routed pieces versus Full F3}

For a fixed-depth routed F3 piece, the present interface distinguishes the high-conductor kernel from the entire
piece.  The needed inputs are:
\[
 \begin{array}{ll}
 (1)&\text{high-conductor kernel power saving},\\
 (2)&\text{transfer of that kernel estimate to the routed source},\\
 (3)&\text{low residual-conductor / conductor-window saving},\\
 (4)&\text{routed main-term cancellation},\\
 (5)&\text{middle-conductor large sieve}.
 \end{array}
\tag{F3-inputs}
\]
Together they give only the complete routed-piece logarithmic estimate
\[
 \boxed{\mathcal P_r\ll_A X(\log X)^{-A}.}
\tag{F3-route}
\]
It is important that a high-conductor kernel power saving alone does \emph{not} imply (F3-route).

The fixed-depth route does not exhaust Full F3.  In the current family decomposition the one-outer-variable
piece is routed, the two-outer-variable kernel remains open, and the three-/four-outer source dictionaries remain
separate obligations.  Thus
\[
 \boxed{\text{fixed-depth routable F3}\not\Rightarrow\text{Full F3}.}
\tag{F3-notfull}
\]

\subsection{The actual Full-Type-II assembly theorem}

Let $F_1,F_2,F_3$ denote the completed versions of the three families above.  The final project theorem has the
logical shape
\[
 \boxed{
 F_1+F_2+F_3
 +\text{source-exhaustive packet census}
 \Longrightarrow \text{uniform project Type II}.}
\tag{F123-assemble}
\]
A separate transference theorem must then identify that uniform project estimate with the Ford--Maynard Type-II
predicate for the actual $w_n$.  Only after both arrows have been proved does one obtain $\FullFM$.


\section{The Full-Type-II theorem that the programme actually needs}

A collection of source-specific estimates is not yet Ford--Maynard Type II.  The theorem required downstream is
universal in divisor-bounded coefficients and covers the full interval.

\begin{hypothesis}[Source-exhaustive Full Type II]\label{hyp:FullFM}
For the actual shifted-prime difference sequence $w_n$ and every fixed $A>0$, sufficiently small fixed
$\varepsilon>0$, and all divisor-bounded coefficient sequences $\xi_m,\kappa_n$,
\[
 \boxed{
 \left|\sum_{\substack{(X/2)^\varepsilon<m\le X^{1/6-\varepsilon}\\X/2<mn\le X}}
 \xi_m\kappa_n w_{mn}\right|
 \ll_A X(\log X)^{-A}.}
\tag{FullFM}
\]
Moreover the global source decomposition is exhaustive: every physical packet is either an exact instance of a
proved packet theorem, a controlled finite-template specialization with total operator cost $X^{o(1)}$, or an
exceptional packet already below target.
\end{hypothesis}

Gate 1A can be one provider for some canonical packets if its five source hypotheses are proved and its scale
is matched.  Gate 1B is intended to provide the complementary signed/high-order packets.  In addition, the
F1/F2/F3 assembly in Section~\ref{sec:F123} and the source-exhaustive packet census must be completed.  Neither
Gate 1A nor Gate 1B, separately or by mere juxtaposition, is the universal theorem above.


\section{Proof-specific Ford--Maynard Type-II landing and the source adapter}
\label{sec:fm-structured}

The universal hypothesis (FullFM) remains the clean published sufficient wrapper.  A proof audit of
Ford--Maynard Proposition~7.22 identifies a narrower proof-specific landing point.  After the finite factor
decomposition, the proof reaches their equation (7.23), a finite-depth product sum with one selected product in
the Type-II interval.  Grouping the factors gives
\[
 Y_1(n')=
 \sum_{n'=\prod_{e\in E}n_e}\prod_{e\in E}x_e(n_e),\qquad
 Y_2(n'')=
 \sum_{n''=\prod_{e\notin E}n_e}\prod_{e\notin E}x_e(n_e),
\tag{FM-Y12}
\]
and the universal Type-II hypothesis is invoked only after the proof has reached
\[
 \boxed{
 \sum_{n=n'n''\asymp x}Y_1(n')Y_2(n'')w_{n'n''}.}
\tag{FM-7.23-gen}
\]
Therefore a theorem controlling every \emph{actual generated instance} of (FM-7.23-gen) is sufficient for this
particular invocation of Type II.

\begin{hypothesis}[FM-SIEVEGEN-TYPEII45]\label{hyp:FMSieveGen}
Every grouped bilinear sum of the form (FM-7.23-gen) actually generated by the chosen rerun of the
Ford--Maynard lower-bound proof satisfies the required uniform estimate.
\end{hypothesis}

This is a proof-specific sufficient target, not a Ford--Maynard theorem statement.  The factor depth is fixed by
the sieve parameters: in the relevant decomposition
\[
 \boxed{N=k\ell+r+s,}
\tag{FM-depth}
\]
with $k,\ell,r,s$ bounded in terms of the fixed sieve data.  Thus an ``unbounded convolution depth'' obstruction is
not present at this stage.

The generated one-variable factors belong to a finite-depth grammar built from M\"obius factors, box/support
cutoffs, Mellin/Perron twists, least/greatest-prime-factor phases, perfect-power pullbacks and bounded-depth
convolutions.  We retain the narrower target
\[
 \boxed{\mathrm{FM\!\mbox{-}PERRON\!\mbox{-}GENERATED\!\mbox{-}TYPEII45}.}
\tag{FM-PerronGen}
\]
The current Ford-side archaeology records the following internal programme passes:
\[
\begin{array}{ll}
\mathrm{FM\!\mbox{-}GENERATED\!\mbox{-}ONEVAR\!\mbox{-}ARCHAEOLOGY45} & \text{pass},\\
\mathrm{FM\!\mbox{-}FINITE\!\mbox{-}NUCLEAR\!\mbox{-}COMPILER45} & \text{pass},\\
\mathrm{FM\!\mbox{-}OUTER\!\mbox{-}COEFFICIENT\!\mbox{-}ABSORPTION45} & \text{pass},\\
\mathrm{WEIGHT\!\mbox{-}DEPENDENCE\!\mbox{-}FINITE\!\mbox{-}COMPILER45} & \text{pass}.
\end{array}
\tag{FM-archaeology}
\]
These are programme-level source/compiler audits, not published Ford--Maynard theorems.

\subsection{The unresolved Gate-output/Ford-source adapter}

The next issue is not another abstract rank-reduction theorem.  It is the literal source dictionary from each
actual Gate packet to the generated one-variable factors entering (FM-7.23-gen).  For every physical packet the adapter must provide all of the following data:
\begin{center}
\fbox{\begin{minipage}{0.86\textwidth}
\begin{itemize}[leftmargin=1.5em,itemsep=1pt,topsep=2pt]
\item operator and variable ranges;
\item coefficient and multiplicity;
\item weight law and centering;
\item gcd and collision restrictions;
\item original prefactor and normalized target;
\item route into Gate 1A, Gate 1B, or a direct exceptional class.
\end{itemize}
\end{minipage}}
\end{center}
We label this source-exact splice
\[
 \boxed{\mathrm{SOURCE\!\mbox{-}ADAPTER45}.}
\tag{SOURCE-ADAPTER45}
\]
Its current status is \emph{primary downstream source interface; source-exhaustive dictionary open}.  In
particular the actual weight law must be identified, not guessed: a packet must be shown to have a common
$W_D$, a finite-template $W_{D,e}$, or a genuinely edge-dependent $W_{D,e}$.

The corresponding coordinate census is
\[
 \boxed{\mathrm{FM\!\mbox{-}TO\!\mbox{-}GATE\!\mbox{-}COORDINATE\!\mbox{-}CENSUS45}:
 \text{ source-blocked}.}
\tag{FM-Gate-census}
\]
Once the literal source census exists, the finite coordinate/reassembly compiler is already available at
programme level.

\begin{hypothesis}[Ford--Maynard replacement certificate $C_{\rm FM}$]\label{hyp:CFM}
The source adapter and coordinate census are complete, every Type-II invocation in the chosen rerun lands in a
grouped coefficient pair controlled by Hypothesis~\ref{hyp:FMSieveGen} (or the certified Perron-generated
subclass), and all changes of variables, truncations, suprema, Mellin/Perron decompositions and nuclear
expansions preserve the required norms uniformly.
\end{hypothesis}

Only
\[
 \boxed{\mathrm{SOURCE\!\mbox{-}ADAPTER45}+C_{\rm FM}
 +\mathrm{FM\!\mbox{-}SIEVEGEN\!\mbox{-}TYPEII45}}
\tag{FM-structured-package}
\]
(or the fully certified Perron-generated variant) can replace the published universal Type-II input.  Until the
literal adapter and replacement certificate are constructed, the universal theorem remains the publication-safe
sufficient route.

\section{Ford--Maynard rough sieve weight}

The published sieve construction uses a vector function $g$ and a weight $H(n)$.  In the central family, the
relevant rough-support identity is the following published lemma.

\begin{theorem}[Ford--Maynard rough-support identity, specialized notation]\label{thm:FM718}
Let $g$ be an admissible vector function and $H$ the sieve weight constructed from it in the Ford--Maynard
Section-7 procedure.  On the relevant set $\mathcal N$,
\[
 \boxed{H(n)=(1\star g)(v(n))\,\1_{\{P^-(n)\ge n^\nu\}},
 \qquad |H(n)|\ll_g1.}
\tag{FM-H}
\]
If $(1\star g)$ has fixed sign on the geometric region, then $H(n)$ has the same sign on $\mathcal N$.
\end{theorem}

For the epsilon-shrunk operating family set
\[
 \sigma:=\frac16-2\varepsilon.
\]
If $0\le\varepsilon\le(1/6)/100$, then
\[
 \sigma\ge\frac{49}{300}.
\]

\begin{lemma}[Rough numbers have at most six prime factors]\label{lem:omega6}
If every prime divisor of $n>1$ is at least $n^\sigma$ and $\sigma\ge49/300$, then
\[
 \boxed{\Omega(n)\le6,}
\]
counting prime factors with multiplicity.
\end{lemma}
\begin{proof}
If $k=\Omega(n)$, then $n\ge n^{k\sigma}$, hence $k\sigma\le1$.  Therefore
$k\le1/\sigma\le300/49<7$.
\end{proof}

\section{Published Ford--Maynard five-collar geometry}
\label{sec:fm-five-collar}

For
\[
 P_\varepsilon=\left(\frac12-\varepsilon,\ \varepsilon,\ \nu-2\varepsilon\right),
\]
the proof of Ford--Maynard Theorem~8.3 places every exceptional vector $x\in H_2$ in a five-collar geometry:
one coordinate has one of the forms
\[
 \boxed{
 \frac12+u,\quad \nu+u,\quad 2\nu+u,\quad 1-\nu+u,\quad 1-2\nu+u,
 \qquad |u|\le6\varepsilon.}
\tag{FM-five-collars}
\]
Ford--Maynard Definition~5.1 uses
\[
 \boxed{v(n)=\left(\frac{\log p_1}{\log n},\ldots,
                     \frac{\log p_k}{\log n}\right),\qquad n=p_1\cdots p_k,}
\tag{FM-vn}
\]
so the collar coordinate corresponds to an individual prime factor.  At $\nu_0=1/6$ the five centres are
\[
 \boxed{\frac12,\ \frac16,\ \frac13,\ \frac56,\ \frac23.}
\tag{FM-collar-centres}
\]
This is a published-source geometric input.  It gives a canonical way to select one collar prime in the rough
exceptional sector and materially sharpens the $N_2$ source decomposition below.

\section{Ford--Maynard Theorem 8.2: where pointwise boundedness is actually used}
\label{sec:FM82-localization}

The bounded theorem is useful here because its proof isolates the pointwise condition to one exceptional term.
Let $\mathcal N_1$ be the rough vectors in the main geometric region $H_1$ and let
$\mathcal N_2=\mathcal N\setminus\mathcal N_1$.  The proof of Ford--Maynard Theorem~8.2 begins with
\[
\begin{aligned}
 \sum_p w_p
 &\ge -\sum_p b_p+\sum_{n\in\mathcal P\cup\mathcal N_1}(w_n+b_n)H(n)\\
 &=\sum_{n\in\mathcal N_1}b_nH(n)
   +\sum_{n\in\mathcal P\cup\mathcal N}w_nH(n)
   -\sum_{n\in\mathcal N_2}w_nH(n).
\end{aligned}
\tag{FM82-decomp}
\]
The first term is evaluated from the comparison condition (b.2) through their Lemma~7.20.  The second is
controlled by the ordinary Type-I/Type-II/growth machinery (Proposition~7.19 and Lemma~7.21), giving
\[
 \sum_{n\in\mathcal P\cup\mathcal N}w_nH(n)\ll \frac{x}{(\log x)^2}.
\tag{FM82-middle}
\]
Only at the final term does the published proof invoke the pointwise bounded-class hypothesis:
\[
 n\in\mathcal N_2:\qquad |w_n|\le\tau(n)^\varrho\ll_{\varrho,\nu}1.
\tag{FM82-pointwise}
\]
It then converts the count of $\mathcal N_2$ into the geometric $H_2$ volume.

Thus the pointwise component of (4.1) is not needed to obtain (FM82-decomp), the comparison main term, or
(FM82-middle); it is used to estimate the last $\mathcal N_2$ contribution.  The other part of (4.1), the
comparison prime-mass lower bound
\[
 \sum_p b_p\ge \frac{x}{\varrho\log x},
\tag{FM82-mass}
\]
is still needed to normalize $x/\log x$-scale errors by $\sum_pb_p$.  For the 287 comparison sequence this
lower bound follows from (G0-b1) for some fixed sufficiently large $\varrho$.

\begin{proposition}[Source-specific replay of Ford--Maynard Theorem 8.2]\label{prop:FM82source}
Fix $P=(\gamma,\theta,\nu)$ and a Ford--Maynard sieve function $g$ with main/exceptional regions $H_1,H_2$.
Assume:
\begin{enumerate}[label=(\roman*)]
\item the ordinary sequence conditions (b.1), (b.2), (I), (II), and (w);
\item for some fixed $\varrho_0\ge1$,
\[
 \sum_p b_p\ge \frac{x}{\varrho_0\log x};
\]
\item the source-specific exceptional estimate
\[
 \left|\sum_{n\in\mathcal N_2}w_nH(n)\right|
 \le C_{N_2}\frac{x}{\log x}
 \sum_k\int_{H_2\cap R_k}\frac{du}{u_1\cdots u_k}
 +o\!\left(\sum_pb_p\right).
\tag{FM82-source-N2}
\]
\end{enumerate}
Then the same proof gives
\[
\begin{aligned}
 \frac{\sum_pa_p}{\sum_pb_p}
 \ge&\ 1+
 \sum_k\int_{H_1\cap R_k}
 \frac{(1\star g)(u)}{u_1\cdots u_k}\,du\\
 &\quad -O_g(1/B)-o(1)
 -C_{N_2}\varrho_0
 \sum_k\int_{H_2\cap R_k}\frac{du}{u_1\cdots u_k}.
\end{aligned}
\tag{FM82-source}
\]
\end{proposition}
\begin{proof}
Insert the exact decomposition (FM82-decomp).  Lemma~7.20 gives the $b_nH(n)$ main integral with relative
error $O_g(1/B)$.  Proposition~7.19 and Lemma~7.21 give (FM82-middle), which is
$o(\sum_pb_p)$ because of the prime-mass lower bound.  Apply (FM82-source-N2) to the final term and use
$x/\log x\le\varrho_0\sum_pb_p$.  Dividing by $\sum_pb_p$ yields (FM82-source).
No pointwise bound $|w_n|\le\tau(n)^\varrho$ is used.
\end{proof}

\section{The source-specific rough $N_2$ replacement via the five collars}
\label{sec:N2-collar}

The common normalization $a_n,b_n\mapsto a_n/\log X,b_n/\log X$ is not used as a bounded-class repair: it also
shrinks the comparison prime mass by an additional logarithm.  Instead we keep the physical source and use the
rough geometry directly.

Set
\[
 \nu_0=\frac16,\qquad \sigma=\nu_0-2\varepsilon.
\]
For $0<\varepsilon\le\nu_0/100$, Lemma~\ref{lem:omega6} gives $\Omega(n)\le6$ on the $\sigma$-rough sector.

\subsection{Squareful rough terms are negligible}

If $p^2\mid n\le X$ and $p\ge X^\sigma$, then the total von-Mangoldt-weighted squareful contribution is bounded by
\[
 \sum_{X^\sigma\le p\le X^{1/2}}
 \left(\frac{X}{p^2}+1\right)\log X
 \ll X^{1-\sigma}\log X+X^{1/2}\log X
 =o\!\left(\frac{X}{\log X}\right).
\tag{N2-squareful}
\]
Thus only the squarefree rough sector contributes at the lower-bound scale.

\subsection{Squarefree sector: collar prime plus largest remaining prime}

By (FM-five-collars), choose canonically a collar prime $q$ whose logarithmic coordinate lies in one of the five
published collars.  Let $r$ be the largest remaining prime and write
\[
 \boxed{n=qdr.}
\tag{N2-qdr}
\]
For the shifted-prime programme source, the active two-linear-form problem is
\[
 \boxed{r,\qquad qdr+2.}
\tag{N2-twin-forms}
\]
For the 287 two-sign specialization the corresponding source forms are
\[
 \boxed{r,\qquad 2qdr-1,\qquad 2qdr+1,}
\tag{N2-287-forms}
\]
with the two signs treated separately by the same upper-sieve mechanism.  The collar variable is narrow, and on
the rough sector the local singular factor is uniformly bounded.

\begin{hypothesis}[Uniform collar two-linear-form sieve]\label{hyp:N2cell}
Uniformly in the five collar choices, the rough prefix $qd$, and the residual interval for the final prime $r$,
the relevant pair of linear forms satisfies the standard upper-bound sieve at the scale
\[
 \sum_{\substack{r\in I(q,d)\\ r\ \mathrm{prime}}}
 \Lambda(qdr+2)
 \ll \frac{|I(q,d)|}{\log X}\,\mathfrak S(q,d)
\tag{N2-twin-sieve}
\]
for the shifted-prime source, and analogously
\[
 \sum_{\substack{r\in I(q,d)\\ r\ \mathrm{prime}}}
 \Lambda(2qdr\pm1)
 \ll \frac{|I(q,d)|}{\log X}\,\mathfrak S_\pm(q,d)
\tag{N2-287-sieve}
\]
for each 287 sign, with the local factors uniformly $O(1)$ on the rough sector.
\end{hypothesis}

The harmonic mass of a prime collar is
\[
 \sum_{q\ \text{in a collar}}\frac1q
 \ll_{\nu_0}\varepsilon+\frac1{\log X}.
\tag{N2-collar-mass}
\]
Consequently the source audit reduces the exceptional shifted-prime lane to
\[
 \boxed{
 \sum_{n\in\mathcal N_2}a_n|H(n)|
 \ll \left(\varepsilon+o(1)\right)\frac{X}{\log X},}
\tag{N2-shifted-bound}
\]
modulo Hypothesis~\ref{hyp:N2cell}.  The comparison side is handled by the published rough-number geometry and
the same five-collar volume.

\begin{hypothesis}[Aggregate $N_2$ replacement]\label{hyp:N2agg}
For the actual difference sequence used in the chosen application, the squareful disposal, collar decomposition,
uniform upper-sieve input and comparison estimate combine to give
\[
 \boxed{
 \left|\sum_{n\in\mathcal N_2}w(n)H(n)\right|
 \le C_{N_2}\frac{X}{\log X}
 \sum_k\int_{H_2\cap R_k}\frac{du}{u_1\cdots u_k}
 +o\!\left(\sum_pb_p\right),}
\tag{N2-agg}
\]
with $C_{N_2}$ uniform on the fixed compatible epsilon range.
\end{hypothesis}

We record the programme label
\[
 \boxed{\mathrm{FM\!\mbox{-}N2\!\mbox{-}SHIFTED\!\mbox{-}PRIME\!\mbox{-}COLLAR\!\mbox{-}SIEVE45}}
\]
with status: \emph{research pass modulo a standard uniform two-linear-form upper sieve and the literal source
dictionary}.  This is a sharper source decomposition than the generic coarse-cell formulation in Draft V9, but
it remains external/uninhabited until the uniform sieve and physical source specialization are supplied.

\section{The epsilon-compatibility gate}
\label{sec:epsilon-compat}

The shrunk parameter is not free.  As $\varepsilon\downarrow0$,
\[
 (X/2)^\varepsilon<m\le X^{1/6-\varepsilon}
\]
widens at \emph{both} ends, so the Type-II theorem becomes harder, even though the Ford--Maynard geometric
lower-bound error tends to zero.  The application must therefore freeze one positive epsilon rather than pass
to a limit inside the arithmetic theorem.

Let $g_{\rm FM}$ denote the fixed sieve function used in the published Ford--Maynard $1/6$ lower-bound
certificate, and write
\[
 \eta_{\rm FM}:=
 1+\sum_k\int_{H(P)\cap R_k}
 \frac{(1\star g_{\rm FM})(u)}{u_1\cdots u_k}\,du>0,
 \qquad P=(1/2,0,1/6).
\tag{FM-margin}
\]
Their central-family geometry gives
\[
 \sum_k\int_{H_2(P_\varepsilon)\cap R_k}\frac{du}{u_1\cdots u_k}
 \ll_{g_{\rm FM}}\varepsilon.
\tag{FM-H2-eps}
\]

\begin{hypothesis}[Compatible fixed epsilon]\label{hyp:epsstar}
There exists one fixed
\[
 \boxed{0<\varepsilon_*\le1/600}
\]
such that:
\begin{enumerate}[label=(\roman*)]
\item the 287 Type-I theorem holds with exponent $1/2-\varepsilon_*$;
\item the source-exhaustive Type-II theorem holds on
$(X/2)^{\varepsilon_*}<m\le X^{1/6-\varepsilon_*}$;
\item the rough $N_2$ replacement is valid uniformly at $\sigma=1/6-2\varepsilon_*$;
\item the geometric loss in Proposition~\ref{prop:FM82source}, including the constant $C_{N_2}\varrho_0$,
is $<\eta_{\rm FM}/2$.
\end{enumerate}
\end{hypothesis}

If the Type-I, Type-II and $N_2$ statements are proved for every sufficiently small fixed epsilon, then
Hypothesis~\ref{hyp:epsstar} follows by choosing one epsilon small enough after the Ford--Maynard constants are
fixed.  No theorem in this dossier takes $\varepsilon\to0$ inside the arithmetic source.


\section{Gate-2 conditional endgame -- corrected source-specific form with two Type-II entrances}

\begin{theorem}[Conditional Gate-2 endgame]\label{thm:Gate2}
Assume:
\begin{enumerate}[label=(\roman*)]
\item the exact Gate-0 comparison conditions (b.1), (b.2), (w) and Type I for the physical source;
\item the rough-support identity and source dictionary needed to use the Ford--Maynard $H(n)$;
\item the source-specific aggregate $N_2$ replacement, Hypothesis~\ref{hyp:N2agg};
\item the compatible fixed epsilon, Hypothesis~\ref{hyp:epsstar};
\item and \emph{one} of the following Type-II packages:
  \begin{enumerate}[label=(\alph*)]
  \item the published sufficient route: source-exhaustive Full Type II, Hypothesis~\ref{hyp:FullFM}; or
  \item the proof-specific route: Hypotheses~\ref{hyp:FMSieveGen} and~\ref{hyp:CFM} (or a certified narrower
  Perron-generated theorem together with the same replacement certificate).
  \end{enumerate}
\end{enumerate}
Then the shifted-prime source has positive prime mass at scale $X$ for all sufficiently large $X$.
\end{theorem}
\begin{proof}
Work at the single fixed point
\[
 P_{\varepsilon_*}=
 \left(\frac12-\varepsilon_*,\ \varepsilon_*,\ \frac16-2\varepsilon_*\right).
\]
Conditions (i)--(iv) provide the ordinary sequence class, the Ford--Maynard rough weight and the direct
source-specific replacement for the final $N_2$ term.  We do not place the unbounded von-Mangoldt source in the
bounded sequence class and we do not apply $C^-_{\rm bd}$ directly to it.

For route (a), Hypothesis~\ref{hyp:FullFM} supplies every Type-II invocation used in the ordinary Ford--Maynard
machinery.  For route (b), the replacement certificate $C_{\rm FM}$ verifies invocation-by-invocation that the
structured theorem supplies the same bilinear estimates.  In either case, Proposition~\ref{prop:FM82source}
replays the Theorem-8.2 decomposition with the pointwise-bounded $N_2$ line replaced by
Hypothesis~\ref{hyp:N2agg}.

The central Ford--Maynard geometry supplies the positive limiting margin $\eta_{\rm FM}$, while
(FM-H2-eps) and Hypothesis~\ref{hyp:N2agg} control the exceptional region.  By
Hypothesis~\ref{hyp:epsstar}(iv), the total geometric/source-specific loss is less than half that margin.
Taking the auxiliary parameter $B$ and then $X$ sufficiently large gives
\[
 \sum_pa_p\ge c_{\rm FM}\sum_pb_p
\]
for some fixed $c_{\rm FM}>0$.  Finally (G0-b1) gives $\sum_pb_p\asymp X/\log X$.
\end{proof}

\begin{status}[Downstream status after the late Gate-1B update]
The published universal Type-II theorem remains a sufficient wrapper.  The structured route is strictly
proof-specific and is \emph{not} activated until the analytic structured estimate, SOURCE-ADAPTER45 / the
FM-to-Gate coordinate census, and the complete replacement certificate $C_{\rm FM}$ are proved.  In both routes the V8 class/shrinkage correction remains
controlling: the unbounded source is handled in the ordinary class and the final pointwise-bounded $N_2$ step is
replaced directly.
\end{status}


\part{Gold closure pivot: repaired log-cofactor lane, fixed certificate, and Kummer backend}

\section{Status firewall for the corrected Gold update}

The fixed-certificate pivot survives, but its first packet census does not simplify as far as Draft~V11
suggested.  The explicit Ford certificate and the transference algebra are retained.  What is withdrawn is the
expectation that the literal leakage census immediately reduces to H8/H9, QK56, Gate~1A and direct packets.
A source-exact smooth-parity packet occurs before that reduction.

\begin{center}
\begin{tabular}{p{0.43\textwidth}p{0.42\textwidth}}
\toprule
Object & Draft-V12 status\\
\midrule
\texttt{LOG-COFACTOR-ASYMPTOTIC287} & proved below; alternative finite route only\\
Current abstract \texttt{LCBeta} & specification incomplete; add $M\le2x$ placement condition\\
\texttt{FIXED-CERTIFICATE287-PIN} & published/source pin PASS; explicit $g_0,g_*,H_*$ written below\\
\texttt{FIXED-CERTIFICATE-TRANSFERENCE287} & proved below with the $\mathcal N_2$ error separated\\
\texttt{FINITE-H8/H9-ONLY-CENSUS} & \textbf{FAIL}; exact counterguard written below\\
Generic seven-prime quadratic-Kummer spectrum & conditional analytic child PASS modulo classical Weil/completion\\
\texttt{H8H9-SOURCE-TO-KUMMER45} & open child/source interface, but not the first fixed-leakage open\\
\texttt{287-FIXED-CERTIFICATE-SMOOTH-PARITY45} & \textbf{FIRST LITERAL ANALYTIC OPEN}\\
\texttt{287-FIXED-CERTIFICATE-LEAKAGE45} & open parent correlation (FCL)\\
Full generated-(7.23) census $+C_{\rm FM}$ & still sufficient; fixed-$g_*$ specialization remains relevant\\
Erd\H{o}s \#287 & open\\
\bottomrule
\end{tabular}
\end{center}

\section{Repairing the log-cofactor supply interface}\label{sec:logcof-asym}

The broad log-cofactor blocker is a useful alternative to the fixed Sophie band, but the abstract supply
predicate must place its adjacent pair inside the actual denominator span.  The minimal repair is to demand
\[
 \boxed{M\le 2x,\qquad x+1\le M.}
 \tag{LC-place}
\]
Indeed Lemma~\ref{lem:half} gives $N\le\lfloor M/2\rfloor\le x$, so $x,x+1\in[N,M]$.
Consequently an abstract predicate which records only $x+1\le M$ is not yet a literal input to an
adjacent-hole theorem.

Let $C(j)$ denote the finite cofactor threshold from the finite obstruction architecture and assume the proved
uniform comparison
\[
 C(j)\le U(j):=j\,j!.
 \tag{LC-C}
\]
For a fixed $0<\eta<1/2$ put
\[
 J(M):=\left\lfloor \eta\frac{\log M}{\log\log M}\right\rfloor.
 \tag{LC-J}
\]

\begin{proposition}[Log-cofactor asymptotic]\label{prop:logcof}
For sufficiently large $M$, if a prime $q$ satisfies
\[
 M\le 2J(M)q,
 \tag{LC-window}
\]
then
\[
 q^2>M,\qquad q>C(2J(M)).
\]
\end{proposition}
\begin{proof}
From (LC-window),
\[
 q\ge \frac{M}{2J(M)}=M^{1-o(1)},
\]
so $q^2>M$ for sufficiently large $M$.  On the other hand, by (LC-C),
\[
 \log C(2J)\le \log(2J)+\log((2J)!)
 \le \log(2J)+2J\log(2J).
\]
Since
\[
 J=(\eta+o(1))\frac{\log M}{\log\log M},\qquad
 \log(2J)=(1+o(1))\log\log M,
\]
we obtain
\[
 \log C(2J)\le (2\eta+o(1))\log M.
\]
Because $2\eta<1$,
\[
 C(2J)\le M^{2\eta+o(1)}<\frac{M}{2J}\le q
\]
for sufficiently large $M$.
\end{proof}

Thus the factorial growth of the finite cofactor threshold is not the remaining difficulty.  The genuine
adjacent supply theorem would still have to produce
\[
 \boxed{\begin{aligned}
 \exists x,q_0,q_1:\quad &M\le2x,\quad x+1\le M,\\
 &q_0\mid x,\quad q_1\mid x+1,\\
 &q_i\text{ prime},\quad M\le2J(M)q_i\quad(i=0,1).
 \end{aligned}}
 \tag{LC-supply}
\]
This simultaneous large-prime-factor correlation remains a deep alternative route; Draft~V12 does not use it as
the controlling analytic theorem.

\section{The explicit Ford certificate and corrected fixed-certificate transference}
\label{sec:fixed-cert}

\subsection{The actual certificate $g_*$}

Ford--Maynard's central construction is not merely existential.  Write
\[
 \nu_0:=0.16623,\qquad P_0:=\left(\frac12,0,\nu_0\right).
\]
Their displayed certificate has
\[
 g_0(\varnothing)=1,\qquad
 g_0(x)=-\mathbf 1_{x\le1/2},
\]
\[
 g_0(x_1,x_2)
 =\mathbf 1_{(1-\nu_0)/2<x_1+x_2<1/2},
\]
\[
 g_0(x_1,x_2,x_3)
 =-\mathbf 1_{x_1+x_2+x_3\le1/2},
\]
and is zero in the remaining dimensions.  The certificate inequality is
\[
 (1\star g_0)(\mathbf x)\le0\qquad(\mathbf x\in\mathcal H(P_0)),
\]
and the associated comparison constant satisfies
\[
 C_0=1+I_3+I_4+I_5+I'_5+I_6\ge 0.000006.
 \tag{FC-1}
\]
These are the published Ford--Maynard certificate data used here.

Fix sufficiently small $\varepsilon_*>0$ and put
\[
 P_*=
 \left(\frac12-\varepsilon_*,\ \varepsilon_*,\ \nu_0-2\varepsilon_*\right),
 \qquad
 \gamma_*:=\frac12-\varepsilon_*,
 \qquad
 \sigma_*:=\nu_0-2\varepsilon_*.
\]
Use the truncated certificate
\[
 g_*(\mathbf x)
 :=
 g_0(\mathbf x)\,
 \mathbf 1_{\{|\mathbf x|\le1/2-2\varepsilon_*\}}.
 \tag{FC-2}
\]
The published shrink analysis gives
\[
 C_*(\varepsilon_*)=C_0+O(\varepsilon_*).
 \tag{FC-3}
\]
Hence there exists at least one fixed $\varepsilon_*>0$ for which
\[
 \boxed{C_*(\varepsilon_*)\ge C_0/2>0.}
 \tag{FC-4}
\]
This is an existence statement unless the implied constant in (FC-3) is made effective.

\begin{status}[Fixed-certificate pin]
The controlling certificate is the $1/2,0,\nu_0$ Ford--Maynard certificate above and its fixed small shrink.
A numerical certificate from an unrelated parameter triple is not interchangeable with this one.
We record \texttt{FIXED-CERTIFICATE287-PIN: PASS} at the published/source level, with effectiveness of a
convenient decimal $\varepsilon_*$ still separate.
\end{status}

\subsection{Literal arithmetic weight $H_*$}

For the canonical split $m=m_1m_2$ satisfying
\[
 P^+(m_1)<n^{\sigma_*}\le P^-(m_2),
\]
define
\[
 G_*(m;n)
 =
 \mathbf 1_{m\le n^{\gamma_*}}\,
 \mu(m_1)\,
 g_*\!\left(\mathbf v(m_2;n)\right),
\qquad
 \boxed{H_*(n):=\sum_{m\mid n}G_*(m;n).}
 \tag{FC-5}
\]
This is the literal fixed-$g_*$ specialization of the Ford--Maynard arithmetic weight.

For a prime $p\asymp X$, only $m=1$ contributes, so
\[
 \boxed{H_*(p)=g_*(\varnothing)=1.}
 \tag{FC-6}
\]
Let $\mathcal N_1$ denote the main composite region where the certificate sign applies, and
$\mathcal N_2$ the exceptional collar.  Then
\[
 H_*(n)\le0\qquad(n\in\mathcal N_1),
 \tag{FC-7}
\]
while $\mathcal N_2$ must be controlled separately by the source-specific $N_2$ estimate; (FC-7) is not
extended to all composites.

\subsection{Corrected leakage decomposition}

Put
\[
 w_X(n)=W(n/X)\,[\Lambda(2n-1)+\Lambda(2n+1)-4B(n)]
\]
and
\[
 \mathcal U_*:=(X/2,X]\setminus
 \bigl(\mathcal P\cup\mathcal N_1\cup\mathcal N_2\bigr).
\]
Define
\[
 \boxed{\mathcal L_*(X):=\sum_{n\in\mathcal U_*}w_X(n)H_*(n).}
 \tag{FC-8}
\]
Also put
\[
 B_X:=\sum_{p\in\mathcal P}b_X(p),\qquad
 \mathcal E_{N_2}(X):=
 \sum_{n\in\mathcal N_2}w_X(n)H_*(n).
\]

\begin{theorem}[Fixed-certificate transference]\label{thm:fixed-cert-transfer}
Assume, for the fixed certificate above,
\[
 \sum_{X/2<n\le X}w_X(n)H_*(n)=o(B_X),
 \tag{FC-9}
\]
\[
 \mathcal L_*(X)=o(B_X),
 \tag{FC-10}
\]
and
\[
 \mathcal E_{N_2}(X)=o(B_X),
 \tag{FC-11}
\]
or more generally that the combined leakage and $N_2$ error are sufficiently small compared with the
positive certificate margin.  Assume also the comparison asymptotic
\[
 \sum_{n\in\mathcal N_1}b_X(n)H_*(n)
 =
 \bigl(C_*(\varepsilon_*)-1+o(1)\bigr)B_X.
 \tag{FC-12}
\]
Then
\[
 \boxed{
 \sum_{p\in\mathcal P}a_X(p)
 \ge
 \bigl(C_*(\varepsilon_*)+o(1)\bigr)B_X.}
 \tag{FC-13}
\]
In particular, if $B_X\gg X/\log X$, the affine prime mass is positive at the required scale.
\end{theorem}

\begin{proof}
By (FC-9)--(FC-11),
\[
 \sum_{n\in\mathcal P\cup\mathcal N_1}w_X(n)H_*(n)=o(B_X).
\]
Since $a_X\ge0$, $H_*(p)=1$ on $\mathcal P$, and $H_*(n)\le0$ on $\mathcal N_1$,
\[
 \sum_{p\in\mathcal P}a_X(p)
 \ge
 \sum_{n\in\mathcal P\cup\mathcal N_1}a_X(n)H_*(n).
\]
Expanding $a_X=b_X+w_X$ and using (FC-12) yields (FC-13).
\end{proof}

For positivity one does not logically need an arbitrary logarithmic saving.  A sufficient quantitative target is
\[
 \boxed{
 |\mathcal L_*(X)|+|\mathcal E_{N_2}(X)|
 \le
 \frac{C_*(\varepsilon_*)}{4}\,B_X
 }
 \tag{FC-14}
\]
for all sufficiently large $X$, together with a correspondingly small total-correlation error.  Because the
available certificate margin is tiny, (FC-14) is still a strong analytic statement, but it is weaker than
demanding $O_A(X(\log X)^{-A})$ for every $A$.

\begin{remark}[What the pivot really saves]
Fixing $g_*$ removes the need to prove a theorem uniformly over all Ford certificates.  It does \emph{not}
make the actual $H_*$ packet family finite-dimensional in the sense needed to reduce immediately to H8/H9.
The next section gives an exact counterguard.
\end{remark}

\section{The first literal fixed-certificate packet: smooth M\"obius parity}
\label{sec:smooth-parity}

The literal fixed-$g_*$ census contains a packet before any H8/H9-only reduction.  In the Ford factorisation
take the cell
\[
 k=0,\qquad J=\varnothing.
\]
Then no prime factor exceeds the rough threshold:
\[
 P^+(n)\le n^{\sigma_*}.
\]
Because $g_*(\varnothing)=1$, the corresponding arithmetic weight is
\[
 H_*(n)
 =
 \sum_{\substack{d\mid n\\ d\le n^{\gamma_*}}}\mu(d)
 =
 \sum_{\substack{d\mid n\\ d\le n^{1/2-\varepsilon_*}}}\mu(d).
 \tag{SP-1}
\]
Thus the fixed leakage contains the exact packet
\[
 \boxed{
 \begin{aligned}
 \mathcal P_{\rm sm}(X)
 :=
 \sum_{\substack{X/2<n\le X\\P^+(n)\le n^{\sigma_*}}}
 &W(n/X)
 \bigl[\Lambda(2n-1)+\Lambda(2n+1)-4B(n)\bigr]\\
 &\times
 \sum_{\substack{d\mid n\\d\le n^{1/2-\varepsilon_*}}}\mu(d).
 \end{aligned}}
 \tag{SP-2}
\]
We denote this exact source packet by
\[
 \boxed{\texttt{287-FIXED-CERTIFICATE-SMOOTH-PARITY45}.}
\]

\subsection{Why the fixed census is not H8/H9-only}

Consider a squarefree balanced $k$-prime cell
\[
 n=p_1\cdots p_k,
 \qquad
 \left|\frac{\log p_i}{\log n}-\frac1k\right|<\eta,
\]
where $1/k+\eta<\sigma_*$.  Choose $\eta$ so small that subset products of at most
\[
 r:=\lfloor k\gamma_*\rfloor
\]
primes are below $n^{\gamma_*}$ and subset products of at least $r+1$ primes are above it.  Then the canonical
split has no large-prime component and
\[
 H_*(n)
 =
 \sum_{j=0}^{r}(-1)^j\binom{k}{j}
 =
 \boxed{(-1)^r\binom{k-1}{r}\ne0.}
 \tag{SP-3}
\]
For small fixed $\varepsilon_*$ this gives, for example,
\[
\begin{array}{c|rrrrrr}
k&7&8&9&10&11&12\\
\hline
H_*(n)&-20&-35&70&126&-252&-462.
\end{array}
\]
Hence H9 genuinely occurs, but so do orders $10,11,12,\ldots$.  Therefore
\[
 \boxed{\texttt{FINITE-H8/H9-ONLY-CENSUS: FAIL}.}
 \tag{SP-4}
\]
Any use of historical ``Full-Nine'' regrouping must prove source-exhaustiveness; the name of the parent operator
does not imply that all larger $\Omega(n)$ cells have disappeared.

\subsection{Parity character of the packet}

At the unshrunk point $\gamma=1/2$, for squarefree $n$ with no exact square-root divisor, pairing
$d\leftrightarrow n/d$ gives
\[
 \mu(n/d)=\mu(n)\mu(d).
\]
When $\mu(n)=1$, the paired full divisor sum forces cancellation below $\sqrt n$; when $\mu(n)=-1$, no analogous
cancellation occurs.  The truncated divisor sum in (SP-1) is therefore a literal divisor-parity detector.
Shrinking from $1/2$ to $1/2-\varepsilon_*$ only introduces a central-divisor collar; it does not remove this
parity mechanism.

The first required analytic estimate on this packet need not be arbitrary-log.  It is enough to obtain a
certificate-compatible constant saving, for example
\[
 \boxed{
 |\mathcal P_{\rm sm}(X)|
 \le c_*\,\frac{X}{\log X},
 \qquad
 c_*<
 \frac{C_*(\varepsilon_*)}{4}
 \times(\text{comparison constant}),
 }
 \tag{SP-5}
\]
together with corresponding bounds for the other fixed-$g_*$ packets.  No provider for (SP-5) is presently
identified in this dossier.

\begin{status}[First literal open]
The exact smooth-parity packet (SP-2) occurs before a reduction to H8/H9 or QK56.  Therefore
\[
 \boxed{\texttt{287-FIXED-CERTIFICATE-SMOOTH-PARITY45}}
\]
is the first literal analytic open exposed by the fixed-certificate census.  The parent FCL correlation remains
open as well.
\end{status}

\section{Generic quadratic-Kummer bilinear estimate}\label{sec:kummer}

Let $p$ be an odd prime, let $\chi_p$ be the quadratic character modulo $p$, and let
$F\in\mathbf F_p[X]$ be a quadratic polynomial with nonzero discriminant.  Put
\[
 K(x):=\chi_p(F(x)).
\]
Let $(\alpha_m)$ be supported on an interval of length $M$ and let $(\beta_n)$ be supported on a positive
interval of length $N<p$ whose elements are nonzero modulo $p$.

\begin{proposition}[Quadratic-Kummer bilinear bound]\label{prop:kummer-bilin}
With the notation above,
\[
 \boxed{
 \left|\sum_{m,n}\alpha_m\beta_nK(mn)\right|^2
 \ll p^{o(1)}\|\alpha\|_2^2\|\beta\|_2^2
 \left[M+\left(\frac{M}{\sqrt p}+\sqrt p\right)N\right].}
 \tag{Kum-1}
\]
The only analytic input is the classical Weil bound for the complete non-square quadratic-character sum,
combined with standard completion for an interval.
\end{proposition}
\begin{proof}
By Cauchy--Schwarz in $m$,
\[
 |S|^2\le\|\alpha\|_2^2
 \sum_m\left|\sum_n\beta_nK(mn)\right|^2.
\]
After expansion, the inner correlation is
\[
 C(n_1,n_2):=\sum_{m\in I}K(mn_1)\overline{K(mn_2)}.
\]
Let $r_1,r_2$ be the two roots of $F$ in an algebraic closure.  The polynomial
$F(n_1X)F(n_2X)$ is a square only when multiplication by $n_2/n_1$ preserves the unordered set
$\{r_1,r_2\}$; this stabilizer has at most two elements.  For these exceptional ratios use the trivial bound
$|C(n_1,n_2)|\le M$.  For all other pairs the complete character sum is non-square, and completion plus Weil
gives
\[
 C(n_1,n_2)\ll p^{o(1)}\left(\frac{M}{\sqrt p}+\sqrt p\right).
\]
The stabilizer pairs contribute $O(M)\|\beta\|_2^2$.  For the remaining pairs use
\[
 \left(\sum_n|\beta_n|\right)^2\le N\|\beta\|_2^2.
\]
Substitution gives (Kum-1).
\end{proof}

Relative to the Cauchy-trivial scale $\sqrt{MN}\,\|\alpha\|_2\|\beta\|_2$, the saving factor is
\[
 p^{o(1)}\left(\frac1N+\frac1{\sqrt p}+\frac{\sqrt p}{M}\right)^{1/2}.
 \tag{Kum-2}
\]
For a seven-prime packet use the labelled $5|2$ split
\[
 m=r_1\cdots r_5\asymp Y^5,\qquad n=r_6r_7\asymp Y^2.
\]
The labelled prime-box multiplicity is $Y^{o(1)}$, so
\[
 \|\alpha\|_2^2\ll Y^{5+o(1)},\qquad
 \|\beta\|_2^2\ll Y^{2+o(1)}.
\]
If
\[
 Y^{5/2-o(1)}\le p\le Y^8,
\]
then (Kum-2) gives
\[
 \left(Y^{-2}+p^{-1/2}+p^{1/2}Y^{-5}\right)^{1/2}
 \le Y^{-1/2+o(1)}=X^{-1/18+o(1)}
\]
for $X=Y^9$.

\begin{status}[Generic Kummer backend: conditional child provider]
The separated generic quadratic-Kummer spectrum has fixed-power capacity $X^{-1/18+o(1)}$ once a physical
packet has been source-exactly reduced to this model.  It is therefore a useful conditional analytic child
provider, not a proof of the fixed-certificate leakage family.  The repeated-root locus and isolated pure
multiplicative-character modes remain in their separate degeneracy routers.
\end{status}

\section{The literal Gate-1B source-to-Kummer interface}\label{sec:source-kummer}

The generic estimate above does \emph{not} by itself close H8 or H9.  The remaining literal theorem is:

\begin{hypothesis}[H8/H9 source-to-Kummer]\label{hyp:source-kummer}
Every actual H8/H9 packet, before source-enlarging inequalities are applied, is decomposed source-exactly into
\[
 \begin{cases}
 \text{a generic separated quadratic-Kummer packet covered by Proposition~\ref{prop:kummer-bilin}},\\
 \text{a principal/zero-mode packet},\\
 \text{a repeated-root degeneracy packet},\\
 \text{an isolated quadratic-character packet routed separately},
 \end{cases}
\]
while preserving the physical prefactor, the $q,p,e$ ranges, coefficient norms, the once-only use of each
M\"obius sign, labelled multiplicity, smooth/Mellin cost and source centering.
\end{hypothesis}

This is the exact interface denoted
\[
 \boxed{\texttt{H8H9-SOURCE-TO-KUMMER45}.}
\]
Thus the generic character estimate is not itself an open spectral wall.  The literal H8/H9 source-to-Kummer
transcription remains an open Gate-1B child interface.  For the fixed-certificate 287 route, however, the
smooth-parity packet in Section~\ref{sec:smooth-parity} occurs earlier and is not known to reduce to this child.

\section{Corrected fixed-certificate 287 route}\label{sec:fixed-leakage}

The fixed-certificate route is logically valid but its packet family is richer than the V11 first census
anticipated.  Define the full leakage parent
\[
 \boxed{
 \mathcal L_*(X)=
 \sum_{n\in\mathcal U_*}
 H_*(n)W(n/X)
 [\Lambda(2n-1)+\Lambda(2n+1)-4B(n)].
 }
 \tag{287-FCL}
\]
We continue to denote the required parent estimate by
\[
 \boxed{\texttt{287-FIXED-CERTIFICATE-LEAKAGE45}.}
\]
The exact smooth-parity child (SP-2) is one literal component of this parent and is presently the first exposed
analytic open.

\begin{status}[Controlling interpretation]
The fixed certificate removes the need for uniformity over arbitrary Ford certificates, but it does not remove
the generated smooth small-prime factors, the truncated M\"obius parity detector, subset-product Type-II
positions, Perron/Mellin factors, or the $u,v$ fragmentations.  Consequently there is not yet evidence that the
fixed-$g_*$ route is analytically much shorter than the fixed-$g_*$ specialization of the generated-(7.23)
route.
\end{status}

The shortest honest dependency graph is therefore
\[
\boxed{
\begin{array}{c}
\text{explicit Ford certificate }g_*+
\text{fixed-certificate transference}\\
\Downarrow\\
\texttt{287-FIXED-CERTIFICATE-SMOOTH-PARITY45}
+\text{other fixed-}g_*\text{ polytopal packets}\\
\Downarrow\\
\begin{cases}
\text{source-exhaustive Full-Nine/QK/Gate regrouping},\\
\text{or a genuinely new direct smooth-parity theorem}
\end{cases}\\
\Downarrow\\
\texttt{287-FIXED-CERTIFICATE-LEAKAGE45}\\
+\ \text{Gate 0 / comparison / }N_2\\
\Downarrow\\
\text{positive affine prime mass}\\
\Downarrow\\
\text{287 finite blocker}\\
+\ \text{effectivity / finite bridge}\\
\Downarrow\\
\text{Erd\H{o}s Problem \#287}.
\end{array}}
\tag{287-fixed-DAG}
\]
The universal Full-Type-II route and the proof-specific generated-(7.23) route remain valid sufficient
alternatives.  They are not mandatory merely as logical wrappers, but the fixed-$g_*$ route has not yet
demonstrated a much smaller analytic core.

\part{The exact 287 splice and what remains}

\section{The universal 287 Full-Type-II specialization (alternative sufficient route)}

A clean publication-safe sufficient theorem for 287 is
\[
 \boxed{
 \sum_{\substack{(X/2)^{\varepsilon_*}<m\le X^{1/6-\varepsilon_*}\\X/2<mn\le X}}
 \xi_m\kappa_n W(mn/X)
 \bigl[\Lambda(2mn-1)+\Lambda(2mn+1)-4B(mn)\bigr]
 \ll_A X(\log X)^{-A}}
\tag{287-II}
\]
for all divisor-bounded $\xi_m,\kappa_n$.  This is $\SophieTI$.  Draft~V12 retains it as a clean sufficient alternative.  The fixed-certificate route
removes the universal-over-certificates wrapper, but its literal smooth-parity census still contains a genuine
parity-breaking family and therefore has not been shown analytically shorter than the relevant fixed-$g_*$
generated specialization.

A future Full-Type-II theorem can be used only after verifying all of the following, explicitly:
\begin{enumerate}
\item the physical affine forms are exactly $2n-1$ and $2n+1$;
\item the comparison term is exactly $4B(n)$ with the local density proved above;
\item the coefficient quantifiers are universal over divisor-bounded sequences;
\item the lower endpoint is $(X/2)^{\varepsilon_*}$ and the upper exponent is $1/6-\varepsilon_*$ at the \emph{same fixed compatible} $\varepsilon_*$ from Hypothesis~\ref{hyp:epsstar};
\item the smooth support $W$ is permitted without changing the theorem class;
\item no M\"obius sign, Gauss factor, spectral saving, or fixed-power margin is used twice.
\end{enumerate}

Gate 1A is not logically mandatory for this splice.  If Full Type II is proved directly from Gate-1B plus other
packets, Gate 1A may remain an auxiliary theorem.  If Gate 1A is used, its scale dictionary and all five source
hypotheses must be supplied.


The universal estimate (287-II) remains the clean publication-safe sufficient socket.  If the proof-specific
Ford--Maynard route in Section~\ref{sec:fm-structured} is completed, one may instead seek a narrower 287 theorem
for only the generated coefficient pairs used by the sieve proof.  That replacement is legitimate only after a
287-specific version of $C_{\rm FM}$ certifies that every grouped coefficient pair produced from
\[
 W(n/X)\,[\Lambda(2n-1)+\Lambda(2n+1)-4B(n)]
\]
belongs to the proved generated grammar.  The pair-modulus Gate-1B theorem (1B-FMPM45) is a candidate provider
for this narrower route, but neither the grammar certificate nor the pair-modulus theorem is currently proved.

\section{From the sieve lower bound to an interior Sophie witness}

Assume Gate 0, the $N_2$ replacement and $\SophieTI$.  Gate 2 yields
\[
 \sum_{q\ \mathrm{prime}}W(q/X)
 \bigl(\Lambda(2q-1)+\Lambda(2q+1)\bigr)
 \gg_W\frac X{\log X}.
\tag{287-mass}
\]
Proper prime powers among $2q\pm1$ contribute only
$O(X^{1/2}(\log X)^{O(1)})=o(X/\log X)$, so (287-mass) forces a genuine prime $q$ in the support for which at
least one of $2q\pm1$ is prime.  Corollary~\ref{cor:landing} then excludes every sufficiently large counterexample.

\section{Effectivity and the finite remainder}

A sufficiently-large contradiction is not yet a complete solution.  One must make effective every threshold
used in:
\begin{itemize}
\item weighted maximal Bombieri--Vinogradov / progression estimates;
\item the comparison-sequence asymptotics;
\item the two-linear-form upper sieve and prefix-volume bounds;
\item the Ford--Maynard lower-bound constant at the chosen $\varepsilon$;
\item the smooth partitions and all $o(1)$ terms in the source compilers.
\end{itemize}
Let $M_0$ be the resulting computable threshold.  Below $M_0$, it is unnecessary to enumerate all subsets.
A single Sophie witness $q$ certifies a whole interval of $M$ satisfying the blocker inequalities.  The finite
certificate can therefore be a list
\[
 (q_i,p_i,[L_i,U_i])
\]
whose intervals cover the required range, with the deeper finite obstruction bank used only on uncovered values.

\section{Full theorem-by-theorem closure ledger}

\begin{longtable}{p{0.27\textwidth}p{0.18\textwidth}p{0.46\textwidth}}
\toprule
Statement & Status & Exact remaining obligation\\
\midrule
\endfirsthead
\toprule
Statement & Status & Exact remaining obligation\\
\midrule
\endhead
Half-range placement & proved here & none\\
Top-layer congruence & proved here & none\\
Isolated-hole lemma & proved here & none\\
$\pm$ Sophie finite blocker & proved here & none\\
Fixed band $M\ge9$ and $q>3$ & proved here & none; strict endpoint in $q^2>M$ checked explicitly\\
Comparison mean $B(dm)$ & proved here & none\\
FM conditions (I),(II),(w),(b.1),(b.2),(4.1) & published definitions & application must satisfy them literally\\
FM $1/6$ central/bounded geometry & published theorem & ordinary unshrunk and bounded shrunk classes kept distinct\\
Gate-0 (b.1) / prime mass & derived here from PNT & make constants/source citation explicit\\
Gate-0 (w) & elementary PASS & $B(n)\ll\log\log(3n)$ gives $\sum|w|\tau\ll X(\log X)^2$\\
Gate-0 (b.2) & standard analytic reduction & write uniform multivariate-PNT proof/reference\\
Gate-0 Type I & source theorem needed & exact weighted maximal BV + comparison progression error\\
Gate-1A determinant / BPP finite mechanism & proved conditionally from displayed inputs & none in the finite algebra\\
Gate-1A PB transport & open source theorem & prove (1A-PB), ceiling and bridge\\
Gate-1A fixed-state exclusion & open source theorem & prove polynomial-height exclusion dictionary\\
Gate-1A outer-root pin & open source theorem & prove physical energy identification (1A-root-pin)\\
Gate-1A projective/diagonal bound & strong open source theorem & prove target-scale diagonal/projective bound (1A-diag)\\
Gate-1A recombination & open source theorem & prove literal (1A-rec)\\
Gate-1A canonical theorem & conditional theorem & follows when five preceding pins are supplied\\
Gate-1B low/medium conductor & research derivation & retain source-normalisation pins\\
Gate-1B source split / spent $\mu(e)$ & exact algebra & none on clean squarefree sector\\
Gate-1B character collapse & proved here & none\\
Gate-1B hybrid Poisson / projector & proved here at displayed algebraic level & literal source conventions must match\\
Gate-1B H7 fixed-power repair & conditional proof & source norm Hypothesis~\ref{hyp:H7norm}\\
Gate-1B historical H7 log target & historical open target & superseded as sole controlling frontier by the late short-short/source split\\
Gate-1B H8 short-short & research capacity / source pin open & same $X^{-1/18}$ capacity audit; robustness/source theorem still required\\
Full-Nine / H9 shifted parent & source-linked shifted child open & recover literal cyclic $TT^*$ source, then prove SHIFT-SOURCE-LINKED-CHAR45\\
QK56 full covariance & internal analytic pass modulo source pins & same-$q$ cross-character margin $X^{-5/36+o(1)}$; cross-$q$ spread transcription pending\\
Gate-0 physical source $a_x,b_{x,z},w_{x,z}$ & written explicitly & source-to-FM normalization / analytic inputs\\
Gate-0 finite local density / Euler product & proved here & infinite-product and asymptotic source pins only where used\\
Gate-0 E-separation / CRT / shift / covariance & proved here & no analytic saving claimed\\
Gate-0/1 fixed determinant and Gauss reassembly & proved here & non-unit/gcd strata require separate routing\\
Gate-1B determinant-2 front algebra & proved here & generic HFMV analytic estimate still separate\\
F1 global centering & proved here & correct physical main-term identification remains source dependent\\
F1 routable migration & conditional & must exhaust all long-M\"obius pieces\\
F2 full family & open & PP/PN/NP/NN plus composite/gcd reassembly; NN tensor open\\
F3 routed fixed-depth sector & conditional / partial & two-outer kernel and remaining source families prevent Full F3\\
Full F1/F2/F3 assembly & open & prove (F123-assemble) and source-exhaustive census\\
Source-exhaustive universal Full Type II & open sufficient wrapper & prove Hypothesis~\ref{hyp:FullFM} plus packet census\\
FM-SIEVEGEN-TYPEII45 & proof-specific open interface & prove Hypothesis~\ref{hyp:FMSieveGen} for every actual generated (7.23) instance\\
FM-PERRON-GENERATED-TYPEII45 & proof-specific sufficient target; grammar/source certificate open & fixed-depth grammar archaeology passed internally\\
SOURCE-ADAPTER45 / FM-to-Gate coordinate census & primary source interface open & print every physical packet, ranges, weight law, multiplicity, prefactor, target and route\\
Ford--Maynard replacement certificate $C_{\rm FM}$ & open finite/source theorem & after the source adapter, certify every generated Type-II invocation and transformation\\
FM rough-support identity & published theorem & exact source instantiation\\
$\Omega(n)\le6$ & proved here & none\\
FM five-collar geometry & published-source pass & no new geometric theorem required\\
$N_2$ squareful rough lane & elementary pass & negligible by (N2-squareful)\\
$N_2$ collar two-form lane & research pass modulo standard upper sieve & prove Hypothesis~\ref{hyp:N2cell} uniformly for the physical source\\
Aggregate $N_2$ replacement & source-specific open splice & instantiate collar sieve + comparison side uniformly at fixed $\varepsilon_*$\\
Source-specific FM-8.2 replay & conditional theorem proved here & ordinary FM machinery + direct $N_2$ replacement; no direct use of $C^-_{\rm bd}$\\
Compatible $\varepsilon_*$ & open compatibility gate & prove Hypothesis~\ref{hyp:epsstar} at one fixed positive epsilon\\
Gate-2 endgame & corrected conditional theorem & Gate 0 + rough $N_2$ + compatible fixed epsilon + either universal Full Type II or certified proof-specific replacement\\
Repaired log-cofactor placement & proved here & abstract supply must include $M\le2x$ as well as $x+1\le M$\\
Log-cofactor threshold asymptotic & proved here & only the simultaneous adjacent large-factor supply remains open on that route\\
Explicit fixed certificate $g_*$ & published/source PASS & effectiveness of a convenient decimal $\varepsilon_*$ remains separate\\
Fixed-certificate transference & proved here & total correlation + leakage + $N_2$ error + positive comparison margin\\
Smooth-parity fixed-certificate packet & FIRST LITERAL ANALYTIC OPEN & prove a certificate-compatible saving for (SP-2)\\
Finite H8/H9-only census & FAIL & higher balanced defect orders occur literally by (SP-3)\\
Generic quadratic-Kummer bilinear backend & conditional child provider & gives $X^{-1/18+o(1)}$ only after source-exact reduction\\
H8/H9 source-to-Kummer & open child/source theorem & relevant to H8/H9 packets but not exhaustive for fixed leakage\\
287 fixed-certificate leakage & open parent correlation & bound all fixed-$g_*$ packets, including smooth parity\\
All-generated-(7.23) census $+C_{\rm FM}$ & alternative sufficient route & fixed-$g_*$ specialization may still be needed to treat smooth small-prime packets\\
287 universal Type-II specialization & open alternative sufficient route & prove (287-II) only if using the universal/full-Type-II path\\
Effectivity & open & compute all thresholds\\
Finite remainder & finite task & interval-cover certificates below $M_0$\\
\bottomrule
\end{longtable}

\section{End-to-end dependency graph}

\begin{center}
\fbox{\begin{minipage}{0.92\textwidth}
\centering
\textbf{Corrected fixed-certificate 287 route}\par
explicit Ford certificate $g_*$
$+$ fixed-certificate transference\par
$\Downarrow$\par
\texttt{287-FIXED-CERTIFICATE-SMOOTH-PARITY45}
$+$ other fixed-$g_*$ packets\par
$\Downarrow$\par
\texttt{287-FIXED-CERTIFICATE-LEAKAGE45}
$+$ Gate 0/comparison/$N_2$\par
$\Downarrow$\par
positive affine prime mass\par
$\Downarrow$\par
interior Sophie witness\par
$\Downarrow$\par
287 finite blocker\par
$\Downarrow$\par
no sufficiently-large counterexample\par
$+$ effectivity $+$ finite certificates\par
$\Downarrow$\par
\textbf{Erd\H{o}s Problem \#287.}
\end{minipage}}
\end{center}

The Gate-1B generic high-order backend now has the separate internal chain
\[
\begin{array}{c}
 \text{actual H8/H9 source}\\
 \Downarrow\ \texttt{H8H9-SOURCE-TO-KUMMER45}\\
 \text{generic quadratic-Kummer packet}\\
 \Downarrow\ \text{Proposition~\ref{prop:kummer-bilin}}\\
 X^{-1/18+o(1)}\text{ saving}.
\end{array}
\]
plus the explicit zero/principal and degeneracy routers.  This is relevant both to the Twin-Prime programme and
to any fixed-certificate leakage decomposition that produces those packets.

\paragraph{Alternative sufficient routes.}
The longer programme remains mathematically valid and important:
\[
\begin{array}{c}
 \text{Gate 0 + Gate 1A + Gate 1B + F1/F2/F3}\\
 +\ \texttt{SOURCE-ADAPTER45}\\[1mm]
 \Downarrow\\[1mm]
 \FullFM\\
 \text{or a certified }C_{\rm FM}+\mathrm{FM\!\mbox{-}SIEVEGEN\!\mbox{-}TYPEII45}.
\end{array}
\]
followed by the corrected rough-$N_2$ Gate-2 replay.  Draft~V12 no longer treats completion of this whole
wrapper as logically mandatory for Problem~\#287 once the fixed-certificate route is available.

\section{What this document claims and does not claim}

The finite unit-fraction part through Corollary~\ref{cor:landing} is written as a complete mathematical proof.
The comparison mean is also proved.  The Ford--Maynard theorem statements are external published inputs.
Gate 0 is written from its physical shifted-prime detector through local densities, centering, CRT, fixed-shift
determinants and the exact Ford--Maynard sequence conditions.  Gate 1A is a conditional theorem whose complete
current mechanism and five load-bearing source hypotheses are displayed.  Gate 1B includes the determinant-2 and HFMV front algebra, the historical H7 dual-determinant repair, the late
H7 short-short $X^{-1/18}$ conditional compiler, the H8 capacity audit, the full $4/9$ complement routing, the
Full-Nine shifted-quotient parent, the QK56 covariance parent, the canonical zero-mode pass and the moving
pair-modulus source-multiplier frontier.  The F1/F2/F3 assembly preceding universal Full Type II is written out
separately, and the later Ford--Maynard proof audit adds a second, proof-specific structured Type-II entrance
which remains conditional on the literal SOURCE-ADAPTER45 packet dictionary and a complete replacement certificate.  Gate 2 is a corrected conditional theorem: the dossier replays the Ford--Maynard Theorem-8.2 decomposition
for the ordinary class and replaces the pointwise-bounded $N_2$ step by an explicit source-specific open lemma.
It does not apply bounded-class positivity directly to the unbounded 287 source.
Accordingly this dossier is designed to make the remaining mathematical work auditable; it is not a proof that
the remaining hypotheses hold.

Draft~V12 additionally pins the explicit Ford certificate, writes its literal arithmetic weight,
proves the corrected fixed-certificate transference identity with the $N_2$ error separated, and exhibits an
exact smooth-parity child of the leakage family.  Balanced squarefree cells prove that the fixed census is not
H8/H9-only.  The generic quadratic-Kummer backend with $X^{-1/18+o(1)}$ capacity is retained only as a
conditional child provider.  The document does not claim the H8/H9 source-to-Kummer transcription, the
smooth-parity estimate, the full fixed-certificate leakage estimate, or Problem~\#287.

\section{Draft V12 corrected Gold-pivot delta}

Relative to Draft~V11, the controlling corrections are:
\begin{enumerate}
\item The fixed Ford certificate is no longer abstract: the dossier records the explicit central $g_0$, the
shrunk fixed $g_*$, the positive certificate constant $C_*(\varepsilon_*)$, and the literal arithmetic weight
$H_*$.
\item The transference theorem is corrected to separate the main composite region $\mathcal N_1$ from the
exceptional collar $\mathcal N_2$.  Total correlation, fixed leakage and the $N_2$ source error are distinct
antecedents.
\item Arbitrary-log leakage saving is no longer presented as logically necessary.  A sufficiently strong
certificate-margin constant saving is enough for positivity.
\item The expectation
\[
 \text{fixed leakage}\longrightarrow
 \{\mathrm{H8/H9},\mathrm{QK56},\mathrm{1A},\mathrm{direct}\}
\]
is withdrawn.  The exact $k=0$, $J=\varnothing$ cell gives the smooth-parity packet (SP-2).
\item Balanced squarefree $k$-prime cells give
\[
 H_*(n)=(-1)^r\binom{k-1}{r}\ne0,
 \qquad r=\lfloor k\gamma_*\rfloor,
\]
so orders $10,11,12,\ldots$ occur literally.  Hence
\texttt{FINITE-H8/H9-ONLY-CENSUS: FAIL}.
\item \texttt{287-FIXED-CERTIFICATE-SMOOTH-PARITY45} is the first literal analytic open exposed by the
fixed-certificate census.  The parent \texttt{287-FIXED-CERTIFICATE-LEAKAGE45} remains open.
\item The seven-prime Kummer estimate remains useful with $X^{-1/18+o(1)}$ capacity, but only as a conditional
analytic provider after source-exact reduction.  It does not close the smooth-parity packet.
\item The fixed-certificate pivot remains logically valid, but the dossier no longer claims evidence that it is
analytically much shorter than the fixed-$g_*$ generated-(7.23) route.
\item All prior safeguards remain in force: the repaired log-cofactor placement, ordinary/bounded
Ford--Maynard firewall, source-specific $N_2$ replacement, fixed compatible $\varepsilon_*$, Gate-1A source
pins, Gate-1B provenance, effectivity and finite completion.
\end{enumerate}

\begin{thebibliography}{18}
\bibitem{ErGr80}
P. Erd\H{o}s and R. L. Graham,
\emph{Old and New Problems and Results in Combinatorial Number Theory},
Monographies de L'Enseignement Math\'ematique, 1980, p.~33.

\bibitem{Va99}
Various authors,
\emph{Some of Paul's Favorite Problems},
booklet for the conference ``Paul Erd\H{o}s and his mathematics'', Budapest, July 1999, Problem~1.15.

\bibitem{FM24}
K. Ford and J. Maynard,
\emph{On the Theory of Prime-Producing Sieves}, arXiv:2407.14368 (2024).

\bibitem{Hux72}
M. N. Huxley,
\emph{On the Difference between Consecutive Primes}, Invent. Math. 15 (1971/72), 164--170.

\bibitem{Zyg02}
A. Zygmund, \emph{Trigonometric Series}, 3rd ed., Cambridge University Press, 2002.

\bibitem{BC18}
S. Bettin and V. Chandee,
\emph{Trilinear Forms with Kloosterman Fractions}, Adv. Math. 328 (2018), 1234--1262.

\bibitem{Wri26}
T. Wright,
\emph{Trilinear Kloosterman Fractions I: Partially Fixed Moduli and Unbalanced Convolutions},
arXiv:2604.25177 (2026).

\bibitem{BP26}
V. Blomer and A. Pascadi,
\emph{Bilinear Forms with Kloosterman Sums via Quadratic Characters},
arXiv:2607.24311 (2026).

\bibitem{Pas25}
A. Pascadi,
\emph{Non-Abelian Amplification and Bilinear Forms with Kloosterman Sums},
arXiv:2511.08445 (2025).
\bibitem{FS22}
E. Fouvry and I. E. Shparlinski,
\emph{On Character Sums with Determinants}, arXiv:2210.15761 (2022).

\bibitem{Shp13}
I. E. Shparlinski,
\emph{Multiple Exponential and Character Sums with Monomials}, arXiv:1309.0324 (2013).

\bibitem{BG12}
J. Bourgain and M. Z. Garaev,
\emph{Sumsets of Reciprocals in Prime Fields and Multilinear Kloosterman Sums}, arXiv:1211.4184 (2012).

\bibitem{FKMS25}
E. Fouvry, E. Kowalski, P. Michel and W. Sawin,
\emph{Bilinear Forms with Trace Functions}, arXiv:2511.09459 (2025).
\end{thebibliography}

\end{document}
